% LaTeX master Datei(en) zusammengestellt von Falk-Jonatan Strube zur Nutzung an der Hochschule für Technik und Wirtschaft Dresden: https://github.com/genericFJS/htw
\documentclass{scrreprt}
\gdef\pathtomaster{../_LaTeX_master}
%% HTW Dresden Corporate Design
%% Copyright (C) Falk-Jonatan Strube <falkjonatanstrube@gmail.com>, 2017
%%% === Anfang Paket htwscr ===
\makeatletter
\ifcsname KOMAoption\endcsname%
\KOMAoption{footheight}{11mm}	% mehr Platz für Wasserzeichen o.ä. in der Fußzeile
\recalctypearea	% s.o.
\RequirePackage{geometry}
\geometry{top=25mm, bottom=25mm, headsep=10mm, inner=25mm, outer=25mm, footskip=15mm, marginparsep=2mm}
\RequirePackage{helvet}	% Corporate Design schreibt Arial/Helvetica vor
\renewcommand{\familydefault}{\sfdefault}	% s.o.
\RequirePackage{graphicx}	% für Grafiken
\RequirePackage{xcolor}		% für Farben
\definecolor{htworange}{RGB}{249,155,28}	% Corporate Design orange CMYK(0/45/100/0)
\definecolor{htwblue}{RGB}{0,116,188}			% Corporate Design blau CMYK(90/50/0/0)
\definecolor{htwgrey}{RGB}{128,128,128}		% Corporate Design gray 60% Schwarz [Annahme: CMYK(0,0,0,60)]
\RequirePackage{ifthen,xifthen}	% für Bedingungen/Optionen
\RequirePackage{etoolbox}				% s.o.

%%% Format des HTW-Headers (wird hier als Texteelement verwenden) nach dem Style der PowerPoint-Vorlagen. Weiterhin inspiriert durch das Paket tudscr (Corporate Design der TU Dresden)
%%% Dateipfad des Logos abhängig von der Verwendung als Paket/Input
\ifdefined\@isHTWCDpackage%
	\gdef\@pathtoimage{HTW}%
\else%
	\ifdefined\pathtomaster%
		\gdef\@pathtoimage{\pathtomaster /HTW}%
	\else 
		\gdef\@pathtoimage{HTW}%
	\fi%
\fi%
\newcommand{\pathtologo}[1]{\gdef\@pathtoimage{#1}}
%%% Berechnung der Größe des Logos
\newlength{\htwLogoHeight}
\setlength{\htwLogoHeight}{12mm}	% Mindesthöhe des Logos laut Corporate Design: 11mm
\newlength{\@htwLogoWidth}
\AtBeginDocument{
\def\printLogo{\includegraphics[height=\htwLogoHeight]{\@pathtoimage}}
\setlength\@htwLogoWidth{\widthof{\printLogo}}
}
%%% Definition des Header
\newcommand{\@htwheader}{{\noindent%
%%% Header ist abhängig von dem Vorhandensein von department, institute und chair
\newcommand{\hasDepInsCh}{\gdef\@hasDepInsCh{}}%
\ifdefined\@department\hasDepInsCh\fi%
\ifdefined\@institute\hasDepInsCh\fi%
\ifdefined\@chair\hasDepInsCh\fi%
\begin{minipage}[b][\htwLogoHeight]{\dimexpr \textwidth - \@htwLogoWidth - 2mm \relax}% 2mm Abstand zwischen Logo und Linie
{\small{\bfseries%
\ifdefined\@faculty\@faculty{}\fi} %
\strut}%
{\color{htworange}\hrule height 1pt}%
\vspace{.1mm}%
\end{minipage}
\hfill%
\begin{minipage}[b][\htwLogoHeight]{\@htwLogoWidth}
\includegraphics[height=\htwLogoHeight]{\@pathtoimage}
\end{minipage}
\ifdefined\@hasDepInsCh%
\\%
\begin{minipage}{\textwidth}
\small%
\ifdefined\@department\@department{}\fi%
\ifdefined\@institute\ifdefined\@department,~\fi \@institute\fi%
\ifdefined\@chair\ifdefined\@department,~\else\ifdefined\@institute,~\fi\fi\@chair\fi\strut
%{\color{htworange}\hrule height 0.5pt}
\end{minipage}
\fi
}}

%%% Definition der Fußzeilen
\usepackage{scrlayer-scrpage}
\@ifclassloaded{scrreprt}{%
\renewcommand*{\partpagestyle}{scrheadings}
\renewcommand*{\chapterpagestyle}{scrheadings}%
}{}
\pagestyle{scrheadings}
\cfoot{\vfill\pagemark}
\chead{}

%%% Einstellungen
\newcommand{\simpledate}{\gdef\@simpledate{}}		% Datum nicht mit "Eingereicht am:" Vorsatz
\newcommand{\nouppercase}{\gdef\@nouppercase{}}	% Titel/Parts/Chapters usw. nicht in Großbuchstaben

%%% Einiges erst zu Beginn des Dokumentes ausführen, um Einstellungen zu berücksichtigen
\AtBeginDocument{
\renewcommand\SS{SS}	% fix big ß
%%% Sektionen usw. ggf. in Großbuchstaben:
\providecommand{\chapterlinesformat}[3]{}	% compatibility
\providecommand{\chapterlineswithprefixformat}[3]{}	% compatibility
\providecommand{\sectionlinesformat}[4]{}	% compatibility
\ifdefined\@nouppercase\relax\else
\renewcommand\sectionlinesformat[4]{\@hangfrom{\hskip#2 #3}{\MakeUppercase{#4}}}
\@ifclassloaded{scrreprt}{%
\renewcommand\chapterlinesformat[3]{\@hangfrom{#2}{\MakeUppercase{#3}}}
\renewcommand\chapterlineswithprefixformat[3]{\MakeUppercase{#2#3}}%
}{}
\fi
%%% Anpassung der Part-Seite
\renewcommand*{\partheadstartvskip}{\@htwheader\vspace*{30mm}}
\renewcommand*{\raggedpart}{\raggedright}
\renewcommand*{\partformat}{\huge\ifdefined\@nouppercase \partname~\thepart\else\MakeUppercase{\partname~\thepart}\fi.}
\ifdefined\@nouppercase\setkomafont{part}{\Huge}\else\setkomafont{part}{\Huge\MakeUppercase}\fi
}

\newcommand{\bilingual}[3]{\ifdefined #1 \let #1\undefined\fi\newcommand{#1}{\iflanguage{english}{#3}{#2}}}	% Übersetzung von einigen Variablen in mehrere Sprachen

\newcommand{\checkThesis}[2]{\ifthenelse{\equal{#1}{diss}\OR\equal{#1}{doctoral}\OR\equal{#1}{phd}}{
\gdef#2{\dissertationname}}{
\ifthenelse{\equal{#1}{diploma}}{
\gdef#2{\diplomathesisname}}{
\ifthenelse{\equal{#1}{master}}{
\gdef#2{\masterthesisname}}{
\ifthenelse{\equal{#1}{bachelor}}{
\gdef#2{\bachelorthesisname}}{
\ifthenelse{\equal{#1}{student}}{
\gdef#2{\studentthesisname}}{
\ifthenelse{\equal{#1}{evidence}}{
\gdef#2{\studentresearchname}}{
\ifthenelse{\equal{#1}{project}}{
\gdef#2{\projectpapername}}{
\ifthenelse{\equal{#1}{seminar}}{
\gdef#2{\seminarpapername}}{
\ifthenelse{\equal{#1}{term}}{
\gdef#2{\termpapername}}{
\ifthenelse{\equal{#1}{research}}{
\gdef#2{\researchname}}{
\ifthenelse{\equal{#1}{log}}{
\gdef#2{\logname}}{
\ifthenelse{\equal{#1}{report}}{
\gdef#2{\reportname}}{
\ifthenelse{\equal{#1}{internship}}{
\gdef#2{\internshipname}}{
\ifthenelse{\equal{#1}{lecture}\OR\equal{#1}{lesson}\OR\equal{#1}{pract}}{
\bilingual{\professorname}{Vorlesung von}{Lecture by}
\simpledate
\gdef\@islecture{\relax}
\ifthenelse{\equal{#1}{lecture}}{\gdef#2{\lecturename}}{
\ifthenelse{\equal{#1}{lesson}}{\gdef#2{\lessonname}}{
\gdef#2{\practname}
}}}{
\gdef#2{#1}}}}}}}}}}}}}}}}

%%% Textfragmente der Titelseite in deutsch und englisch
\bilingual{\dateofbirthname}{Geboren am:}{Born on:}
\bilingual{\placeofbirthname}{in}{in}
\bilingual{\coursename}{Studiengang:}{Course:}
\bilingual{\disciplinename}{Studienrichtung:}{Discipline:}
\bilingual{\matriculationnumbername}{Matrikelnummer:}{Matriculation number:}
\bilingual{\matriculationyearname}{Immatrikulationsjahr}{Matriculation year:}
\bilingual{\graduationname}{zur Erlangung des akademischen Grades}{to achieve the academic degree}
\bilingual{\refereename}{Gutachter}{Referee}
\bilingual{\advisorname}{Fachreferent}{Advisor}
\bilingual{\supervisorname}{Betreuer}{Supervisor}
\bilingual{\professorname}{Betreuender Hochschullehrer}{Supervising professor}
\bilingual{\datename}{Eingereicht am:}{Submitted on:}
%% Thesis Typen (inspiriert von dem Paket tudscr):
\bilingual{\dissertationname}{Dissertation}{Dissertation} % diss, doctoral, phd 
\bilingual{\diplomathesisname}{Diplomarbeit}{Diploma Thesis} % diploma 
\bilingual{\masterthesisname}{Master-Arbeit}{Master Thesis} % master 
\bilingual{\bachelorthesisname}{Bachelor-Arbeit}{Bachelor Thesis} % bachelor 
\bilingual{\studentthesisname}{Studienarbeit}{Student Thesis} % student 
\bilingual{\studentresearchname}{Großer Beleg}{Student Research Project} % evidence 
\bilingual{\projectpapername}{Projektarbeit}{Project Paper} % project
\bilingual{\seminarpapername}{Seminararbeit}{Seminar Paper} % seminar
\bilingual{\termpapername}{Hausarbeit}{Term Paper} % term
\bilingual{\researchname}{Forschungsbericht}{Research Report} % research 
\bilingual{\logname}{Protokoll}{Log} % log 
\bilingual{\reportname}{Bericht}{Report} % report
\bilingual{\internshipname}{Praktikumsbericht}{Internship Report} % internship 
\bilingual{\lecturename}{Vorlesungsmitschrift}{Lecture Notes} % lecture 
\bilingual{\lectureauthorname}{Mitschrift von}{Notes by} % lecture  
\bilingual{\lessonname}{\"Ubungsmitschrift}{Lesson Notes} % lecture  
\bilingual{\practname}{Praktikumsmitschrift}{Pracital Course Notes} % lecture 
%%% Variablen für die Titelseite (inspiriert von dem Paket tudscr):
\renewcommand{\subject}[1]{\checkThesis{#1}{\@subject}}
\newcommand{\faculty}[1]{\gdef\@faculty{#1}}
\newcommand{\department}[1]{\gdef\@department{#1}}
\newcommand{\institute}[1]{\gdef\@institute{#1}}
\newcommand{\chair}[1]{\gdef\@chair{#1}}
\newcommand{\authormore}[1]{\gdef\@authormore{#1}}
\newcommand{\dateofbirth}[1]{\gdef\@dateofbirth{#1}}
\newcommand{\placeofbirth}[1]{\gdef\@placeofbirth{#1}}
\newcommand{\course}[1]{\gdef\@course{#1}}
\newcommand{\discipline}[1]{\gdef\@discipline{#1}}
\newcommand{\matriculationnumber}[1]{\gdef\@matriculationnumber{#1}}
\newcommand{\matriculationyear}[1]{\gdef\@matriculationyear{#1}}
\newcommand{\thesis}[1]{\checkThesis{#1}{\@thesis}}
\newcommand{\graduation}[2][]{\gdef\@graduation{#2}\gdef\@graduationshort{#1}}
\newcommand{\supervisor}[1]{\gdef\@supervisor{#1}}
\newcommand{\referee}[1]{\gdef\@referee{#1}}
\newcommand{\advisor}[1]{\gdef\@advisor{#1}}
\newcommand{\professor}[1]{\gdef\@professor{#1}}
%%% Titelseite (inspiriert von dem Paket tudscr):
\renewcommand*{\maketitle}{%
\begin{titlepage}%
\begin{flushleft}
\renewcommand{\thefootnote}{$\star$}%
%%% Kopfzeile HTW-Style:
\@htwheader\\[13mm]
%%% Die Zeilen für titlehead, subject, title, subtitle und author sind immer vorhanden (spacing)
{\noindent\@titlehead{}\strut}\\[3mm]%
{\Large\@subject{}\strut}\\[8mm]%
{\Huge\bfseries%
\ifdefined\@nouppercase\@title{}\else\MakeUppercase{\@title{}}\fi%
\strut}\\
{\Large\bfseries%
\ifdefined\@nouppercase\@subtitle{}\else\MakeUppercase{\@subtitle{}}\fi%
\strut}\\[11mm]
\ifdefined\@islecture\strut\else{\Large\@author{}\strut}\fi % Autor wird bei Vorlesungsmitschrift unten angezeigt
%%% Weitere Felder werden nur angezeigt, wenn sie definiert wurden
\ifdefined\@authormore\ifdefined\@islecture\else\\\fi{\@authormore{}\strut}\fi
\ifdefined\@dateofbirth\\\dateofbirthname{} \@dateofbirth{} \ifdefined\@placeofbirth\placeofbirthname{} \@placeofbirth{}\strut\fi\fi
\ifdefined\@course\\\coursename{} \@course{}\strut\fi
\ifdefined\@discipline\\\disciplinename{} \@discipline{}\strut\fi
\ifdefined\@matriculationnumber\\\matriculationnumbername{} \@matriculationnumber{}\strut\fi
\ifdefined\@matriculationyear\\\matriculationyearname{} \@matriculationyear{}\strut\fi
\ifdefined\@thesis\\[10mm]{\LARGE\bfseries
\ifdefined\@nouppercase\@thesis{}\else\MakeUppercase{\@thesis{}}\fi
\strut}
\ifdefined\@graduation\\[2mm]{\graduationname{}\strut}
\\[5mm]{\LARGE\bfseries
\ifdefined\@nouppercase\relax\else\MakeUppercase{\fi%
\@graduation{} \ifdefined\@graduationshort(\@graduationshort{})\fi
\ifdefined\@nouppercase\relax\else}\fi\strut}\fi\fi
\vfill\noindent
\ifdefined\@islecture\\[1mm]{\lectureauthorname{}\strut}\\[.7mm]{\Large\@author{}\strut}\fi
\ifdefined\@referee{\refereename{}\strut}\\[.7mm]{\Large\@referee{}\strut}\fi
\ifdefined\@advisor\\[1mm]{\advisorname{}\strut}\\[.7mm]{\Large\@advisor{}\strut}\fi
\ifdefined\@supervisor\\[1mm]{\supervisorname{}\strut}\\[.7mm]{\Large\@supervisor{}\strut}\fi
\ifdefined\@professor\\[1mm]{\professorname{}\strut}\\[.7mm]{\Large\@professor{}\strut}\fi
\ifdefempty{\@date}{}{\ifdefined\@simpledate\\[6mm]\relax\else\\[3mm]\datename{}~\fi\@date{}\strut}
\ifdefempty{\@publishers}{}{\\[4.5mm]{\Large\@publishers{}\strut}}%
\renewcommand\footnoterule{}
\@thanks\global\let\@thanks\@empty
\end{flushleft}
\end{titlepage}
\setcounter{footnote}{0}
}
\else%
\PackageError{htwcd}{htwcd needs a KOMA class!}%
\fi%
\makeatother
%%% === Ende Paket htwscr ===
%-- Nummerierung bis Subsubsection für Report
\setcounter{secnumdepth}{3}	% zählt auch subsubsection
\setcounter{tocdepth}{3}	% Inhaltsverzeichnis bis in subsubsection

% LuaLaTeX Unterscheidung
\usepackage{ifluatex}

% Input inkl. Umlaute, Silbentrennung
\ifluatex
	\usepackage{fontspec}
\else
	\usepackage[utf8]{inputenc}
	\usepackage[T1]{fontenc}
\fi
\usepackage[english,ngerman]{babel}
\usepackage{csquotes}	% Anführungszeichen
\usepackage{marvosym}
\usepackage{eurosym}

% Style-Aufhübschung
\usepackage{soul, color}	% Kapitälchen, Unterstrichen, Durchgestrichen usw. im Text
%\usepackage{titleref}
\usepackage[breakwords, fit]{truncate}	% Abschneiden von Sätzen
\renewcommand{\TruncateMarker}{\,…}

% Mathe usw.
\usepackage{amssymb}
\usepackage{amsthm}
\usepackage[fleqn,intlimits]{amsmath}	% fleqn: align-Umgebung rechtsbündig; intlimits: Integralgrenzen immer ober-/unterhalb
%\usepackage{mathtools} % u.a. schönere underbraces
\usepackage{xcolor}
\usepackage{esint}	% Schönere Integrale, \oiint vorhanden
\everymath=\expandafter{\the\everymath\displaystyle}	% Mathe Inhalte werden weniger verkleinert
\usepackage{wasysym}	% mehr Symbole, bspw \lightning
%\renewcommand{\int}{\int\limits}
%\usepackage{xfrac}	% mehr fracs: sfrac{}{}
\let\oldemptyset\emptyset	% schöneres emptyset
\let\emptyset\varnothing
%\RequirePackage{mathabx}	% mehr Symbole
\mathchardef\mhyphen="2D	% Hyphen in Math
\usepackage{mathtools}		% zum Beispiel {pmatrix}[r] für Formatierung der Matrix

% tikz usw.
\usepackage{tikz}
\usepackage{pgfplots}
\pgfplotsset{compat=1.11}	% Umgeht Fehlermeldung
\usetikzlibrary{graphs}
%\usetikzlibrary{through}	% ???
\usetikzlibrary{arrows}
\usetikzlibrary{arrows.meta}	% Pfeile verändern / vergrößern: \draw[-{>[scale=1.5]}] (-3,5) -> (-3,3);
\usetikzlibrary{automata,positioning} % Zeilenumbruch im Node node[align=center] {Text\\nächste Zeile} automata für Graphen
\usetikzlibrary{matrix}
\usetikzlibrary{patterns}	% Schraffierte Füllung
\usetikzlibrary{shapes.geometric}	% Polygon usw.
\tikzstyle{reverseclip}=[insert path={	% Inverser Clip \clip
	(current page.north east) --
	(current page.south east) --
	(current page.south west) --
	(current page.north west) --
	(current page.north east)}
% Nutzen: 
%\begin{tikzpicture}[remember picture]
%\begin{scope}
%\begin{pgfinterruptboundingbox}
%\draw [clip] DIE FLÄCHE, IN DER OBJEKT NICHT ERSCHEINEN SOLL [reverseclip];
%\end{pgfinterruptboundingbox}
%\draw DAS OBJEKT;
%\end{scope}
%\end{tikzpicture}
]	% Achtung: dafür muss doppelt kompliert werden!
\usepackage{graphpap}	% Grid für Graphen
\tikzset{every state/.style={inner sep=2pt, minimum size=2em}}
\usetikzlibrary{mindmap, backgrounds}
%\usepackage{tikz-uml}	% braucht Dateien: http://perso.ensta-paristech.fr/~kielbasi/tikzuml/

% Tabular
\usepackage{longtable}	% Große Tabellen über mehrere Seiten
\usepackage{multirow}	% Multirow/-column: \multirow{2[Anzahl der Zeilen]}{*[Format]}{Test[Inhalt]} oder \multicolumn{7[Anzahl der Reihen]}{|c|[Format]}{Test2[Inhalt]}
\renewcommand{\arraystretch}{1.3} % Tabellenlinien nicht zu dicht
\usepackage{colortbl}
\arrayrulecolor{gray}	% heller Tabellenlinien
\usepackage{array}	% für folgende 3 Zeilen (für Spalten fester breite mit entsprechender Ausrichtung):
\newcolumntype{L}[1]{>{\raggedright\let\newline\\\arraybackslash\hspace{0pt}}m{\dimexpr#1\columnwidth-2\tabcolsep-1.5\arrayrulewidth}}
\newcolumntype{C}[1]{>{\centering\let\newline\\\arraybackslash\hspace{0pt}}m{\dimexpr#1\columnwidth-2\tabcolsep-1.5\arrayrulewidth}}
\newcolumntype{R}[1]{>{\raggedleft\let\newline\\\arraybackslash\hspace{0pt}}m{\dimexpr#1\columnwidth-2\tabcolsep-1.5\arrayrulewidth}}
\usepackage{booktabs}	% Kompatibilität für array und \hrule
\usepackage{caption}	% Um auch unbeschriftete Captions mit \caption* zu machen

% Nützliches
\usepackage{verbatim}	% u.a. zum auskommentieren via \begin{comment} \end{comment}
\usepackage{tabto}	% Tabs: /tab zum nächsten Tab oder /tabto{.5 \CurrentLineWidth} zur Stelle in der Linie
\NumTabs{6}	% Anzahl von Tabs pro Zeile zum springen
\usepackage{listings} % Source-Code mit Tabs
\usepackage{lstautogobble} 
\usepackage{enumitem}	% Anpassung der enumerates
%\setlist[enumerate,1]{label=(\arabic*)}	% global andere Enum-Items
\renewcommand{\labelitemiii}{$\scriptscriptstyle ^\blacklozenge$} % global andere 3. Item-Aufzählungszeichen
\usepackage{letltxmacro} % neue Definiton von Grundbefehlen
% Nutzen:
%\LetLtxMacro{\oldemph}{\emph}
%\renewcommand{\emph}[1]{\oldemph{#1}}
\RequirePackage{xpatch}	% ua. Konkatenieren von Strings/Variablen (etoolbox)
\usepackage{xstring}	% String Operationen
\usepackage{minibox}	% Minibox anstatt \fbox{} für Boxen mit (manuellen) Zeilenumbruch
\usepackage{framed}		% \begin{framed} statt minibox/fbox mit automatischem Zeilenumbruch
\usepackage{microtype}	% Verfeinerte Hyphenation
\usepackage{import}		
\usepackage{standalone}	% Zusammen mit import: Einfügen von separaten Latex-Dateien (die schon header haben)

% Einrichtung von lst
%\renewcommand*{\ttdefault}{qcr}
\lstset{
basicstyle=\ttfamily, 
%basicstyle=\footnotesize\sffamily\color{black},
%mathescape=true, 
%escapeinside=^^, 
autogobble, 
tabsize=2,
frame=single,
rulecolor=\color{lightgray},
numbers=left,
numbersep=5pt,
numberstyle=\tiny\color{gray},
commentstyle=\color{gray},
keywordstyle=\color{htwblue},
stringstyle=\color{htworange},
morecomment=[l][\color{magenta}]{\#}
showspaces=false,
showstringspaces=false,
breaklines=true,
literate=%
    {Ö}{{\"O}}1
    {Ä}{{\"A}}1
    {Ü}{{\"U}}1
    {ß}{{\ss}}1
    {ü}{{\"u}}1
    {ä}{{\"a}}1
    {ö}{{\"o}}1
    {„}{{\glqq}}1
    {“}{{\grqq}}1
    {~}{{\textasciitilde}}1
}
\usepackage{scrhack} % Fehler umgehen
\def\ContinueLineNumber{\lstset{firstnumber=last}} % vor lstlisting. Zum wechsel zum nicht-kontinuierlichen muss wieder \StartLineAt1 eingegeben werden
\def\StartLineAt#1{\lstset{firstnumber=#1}} % vor lstlisting \StartLineAt30 eingeben, um bei Zeile 30 zu starten
\let\numberLineAt\StartLineAt

% BibTeX
\usepackage[
backend=bibtex,
bibencoding=ascii,
style=authoryear, citestyle=authoryear-ibid,
%doi=false,
%isbn=false,
%url=false
]{biblatex}	% BibTeX
\usepackage{makeidx}
%\makeglossary
%\makeindex
\usepackage{url}	% Darstellung von URLs

% Grafiken
\usepackage{graphicx}
\usepackage{epstopdf}	% eps-Vektorgrafiken einfügen
\usepackage{transparent}	% transparent nutzen: {\transparent{0.4} ...}
%\epstopdfsetup{outdir=\customDir}
% Prüft, ob Grafik existiert (mit \ifvalidimage{}{}) [Quelle: https://tex.stackexchange.com/a/99176]:
\makeatletter
\newif\ifgraphicexist
\catcode`\*=11
\newcommand\ifvalidimage[1]{%
    \begingroup
    \global\graphicexisttrue
    \let\input@path\Ginput@path
    \filename@parse{#1}%
    \ifx\filename@ext\relax
    \@for\Gin@temp:=\Gin@extensions\do{%
        \ifx\Gin@ext\relax
        \Gin@getbase\Gin@temp
        \fi}%
    \else
    \Gin@getbase{\Gin@sepdefault\filename@ext}%
    \ifx\Gin@ext\relax
    \global\graphicexistfalse
    \def\Gin@base{\filename@area\filename@base}%
    \edef\Gin@ext{\Gin@sepdefault\filename@ext}%
    \fi
    \fi
    \ifx\Gin@ext\relax
    \global\graphicexistfalse
    \else 
    \@ifundefined{Gin@rule@\Gin@ext}%
    {\global\graphicexistfalse}%
    {}%
    \fi  
    \ifx\Gin@ext\relax 
    \gdef\imageextension{unknown}%
    \else
    \xdef\imageextension{\Gin@ext}%
    \fi 
    \endgroup 
    \ifgraphicexist
    \expandafter \@firstoftwo
    \else
    \expandafter \@secondoftwo
    \fi 
} 
\catcode`\*=12
\makeatother
\LetLtxMacro{\forceincludegraphics}{\includegraphics}	% neuer Befehl für includegraphics
\renewcommand{\includegraphics}[2][]{	% altes includegraphics neu definieren, damit es auch nicht vorhandene einfügt
\ifvalidimage{#2}{
\forceincludegraphics[#1]{#2}
}{
\message{Achtung: Grafik wurde nicht gefunden: '#2'}
\minibox[frame]{
\textbf{\StrSubstitute{#2}{_}{\_}}%  \ifthenelse{\isempty{#1}}{}{\\\textit{#1}}
}}}

% pdf-Setup
\usepackage{pdfpages}
\usepackage[bookmarks,%
bookmarksopen=false,% Klappt die Bookmarks in Acrobat aus
colorlinks=true,%
linkcolor=black,%
citecolor=black,%
urlcolor=htworange,%
]{hyperref}

% Nach Dokumentenbeginn ausführen:
\makeatletter
\AtBeginDocument{	
	% pdf Informationen
	\hypersetup{
		pdfauthor={\@author},
		pdftitle={\@title},
		pdfproducer={LaTeX Layout erstellt von Falk-Jonatan Strube}
	}
	% signierte Fußzeile
	%\providecommand{\ofootname}{Mitschrift von}
	\providecommand{\pathtomaster}{../}
	\ofoot{%
	%\ifdefined\ofoottext
	%\else
	%\begin{minipage}{40mm}
	%\begin{flushright}
	%\textcolor{gray}{
	%\ofootname\\
	%Falk-Jonatan Strube\vspace*{-1mm}}
	%\end{flushright}
	%\end{minipage}
	%\fi%
	%\hspace*{1mm}%
	\begin{minipage}{10 mm}
	\includegraphics[height=9mm]{\pathtomaster/FJS-logo_gray}
	\end{minipage}%
	}
}
\makeatother



%%% ==================================================
%%% ===================== MACROS =====================
%%% ==================================================
% Horizontale Linie:
\newcommand{\HRule}[1][\medskipamount]{\par
  \vspace*{\dimexpr-\parskip-\baselineskip+#1}
  \noindent\rule[0.2ex]{\linewidth}{0.2mm}\par
  \vspace*{\dimexpr-\parskip-.5\baselineskip+#1}}
% Gestrichelte horizontale Linie:
\RequirePackage{dashrule}
\newcommand{\HDRule}[1][\medskipamount]{\par
  \vspace*{\dimexpr-\parskip-\baselineskip+#1}
  \noindent\hdashrule[0.2ex]{\linewidth}{0.2mm}{1mm} \par
  \vspace*{\dimexpr-\parskip-.5\baselineskip+#1}}
% Mathe in Anführungszeichen:
\newsavebox{\mathbox}\newsavebox{\mathquote}
\makeatletter
\newcommand{\mq}[1]{% \mathquotes{<stuff>}
  \savebox{\mathquote}{\text{"}}% Save quotes
  \savebox{\mathbox}{$\displaystyle #1$}% Save <stuff>
  \raisebox{\dimexpr\ht\mathbox-\ht\mathquote\relax}{"}#1\raisebox{\dimexpr\ht\mathbox-\ht\mathquote\relax}{''}
}
\makeatother

% Paragraph mit Zähler (Section-Weise)
\newcounter{cparagraphC}
\newcommand{\cparagraph}[1]{
\stepcounter{cparagraphC}
\paragraph{\thesection{}-\thecparagraphC{} #1}
%\addcontentsline{toc}{subsubsection}{\thesection{}-\thecparagraphC{} #1}
\label{\thesection-\thecparagraphC}
}
\makeatletter
\@addtoreset{cparagraphC}{section}
\makeatother


% (Vorlesungs-)Folien einbinden:
\makeatletter
% Folien-Pfad (inkl. Dateiname ohne Endung und ggf. ohne Nummerierung)
\newcommand{\@slidePath}{}
\newcommand{\slidePath}[1]{\renewcommand{\@slidePath}{#1}}
% Folien Eigenschaften
\newcommand{\@slideScale}{0.5}
\newcommand{\slideScale}[1]{\renewcommand{\@slideScale}{#1}}
\AtBeginDocument{
% Folien von einer Datei skaliert
\newcommand{\slide}[2][\@slideScale]{\slides[#1]{}{#2}}
\newcommand{\slideTrim}[6][\@slideScale]{\slides[#1 , clip,  trim = #5cm #4cm #6cm #3cm]{}{#2}}
% Folien von mehreren nummerierten Dateien skaliert
\newcommand{\slides}[3][\@slideScale]{\begin{center}
\includegraphics[page=#3, scale=#1]{\@slidePath #2.pdf}
\end{center}}
}
\makeatother

% \emph{} anders definieren
\makeatletter
\DeclareRobustCommand{\em}{%
  \@nomath\em \if b\expandafter\@car\f@series\@nil
  \normalfont \else \scshape \fi}
\makeatother

% unwichtiges
\newcommand{\unimptnt}[1]{{\transparent{0.5}#1}}

% alph. enumerate
\newenvironment{anumerate}{\begin{enumerate}[label=(\alph*)]}{\end{enumerate}} % Alphabetische Aufzählung

% Hanging parameters
\newcommand{\hangpara}[1]{\par\noindent\hangindent+2em\hangafter=1 #1\par\noindent}

%% EINFACHE BEFEHLE

% Abkürzungen Mathe
\newcommand{\EE}{\mathbb{E}}
\newcommand{\QQ}{\mathbb{Q}}
\newcommand{\RR}{\mathbb{R}}
\newcommand{\CC}{\mathbb{C}}
\newcommand{\NN}{\mathbb{N}}
\newcommand{\ZZ}{\mathbb{Z}}
\newcommand{\PP}{\mathbb{P}}
\renewcommand{\SS}{\mathbb{S}}
\newcommand{\cA}{\mathcal{A}}
\newcommand{\cB}{\mathcal{B}}
\newcommand{\cC}{\mathcal{C}}
\newcommand{\cD}{\mathcal{D}}
\newcommand{\cE}{\mathcal{E}}
\newcommand{\cF}{\mathcal{F}}
\newcommand{\cG}{\mathcal{G}}
\newcommand{\cH}{\mathcal{H}}
\newcommand{\cI}{\mathcal{I}}
\newcommand{\cJ}{\mathcal{J}}
\newcommand{\cM}{\mathcal{M}}
\newcommand{\cN}{\mathcal{N}}
\newcommand{\cP}{\mathcal{P}}
\newcommand{\cR}{\mathcal{R}}
\newcommand{\cS}{\mathcal{S}}
\newcommand{\cZ}{\mathcal{Z}}
\newcommand{\cL}{\mathcal{L}}
\newcommand{\cT}{\mathcal{T}}
\newcommand{\cU}{\mathcal{U}}
\newcommand{\cX}{\mathcal{X}}
\newcommand{\cV}{\mathcal{V}}
\renewcommand{\phi}{\varphi}
\renewcommand{\epsilon}{\varepsilon}
\renewcommand{\theta}{\vartheta}

% Verschiedene als Mathe-Operatoren
\DeclareMathOperator{\arccot}{arccot}
\DeclareMathOperator{\arccosh}{arccosh}
\DeclareMathOperator{\arcsinh}{arcsinh}
\DeclareMathOperator{\arctanh}{arctanh}
\DeclareMathOperator{\arccoth}{arccoth} 
\DeclareMathOperator{\var}{Var} % Varianz 
\DeclareMathOperator{\cov}{Cov} % Co-Varianz 

% Farbdefinitionen
\definecolor{red}{RGB}{180,0,0}
\definecolor{green}{RGB}{75,160,0}
\definecolor{blue}{RGB}{0,75,200}
\definecolor{orange}{RGB}{255,128,0}
\definecolor{yellow}{RGB}{255,245,0}
\definecolor{purple}{RGB}{75,0,160}
\definecolor{cyan}{RGB}{0,160,160}
\definecolor{brown}{RGB}{120,60,10}

\definecolor{itteny}{RGB}{244,229,0}
\definecolor{ittenyo}{RGB}{253,198,11}
\definecolor{itteno}{RGB}{241,142,28}
\definecolor{ittenor}{RGB}{234,98,31}
\definecolor{ittenr}{RGB}{227,35,34}
\definecolor{ittenrp}{RGB}{196,3,125}
\definecolor{ittenp}{RGB}{109,57,139}
\definecolor{ittenpb}{RGB}{68,78,153}
\definecolor{ittenb}{RGB}{42,113,176}
\definecolor{ittenbg}{RGB}{6,150,187}
\definecolor{itteng}{RGB}{0,142,91}
\definecolor{ittengy}{RGB}{140,187,38}

% Textfarbe ändern
\newcommand{\tred}[1]{\textcolor{red}{#1}}
\newcommand{\tgreen}[1]{\textcolor{green}{#1}}
\newcommand{\tblue}[1]{\textcolor{blue}{#1}}
\newcommand{\torange}[1]{\textcolor{orange}{#1}}
\newcommand{\tyellow}[1]{\textcolor{yellow}{#1}}
\newcommand{\tpurple}[1]{\textcolor{purple}{#1}}
\newcommand{\tcyan}[1]{\textcolor{cyan}{#1}}
\newcommand{\tbrown}[1]{\textcolor{brown}{#1}}

% Umstellen der Tabellen Definition
\newcommand{\mpb}[1][.3]{\begin{minipage}{#1\textwidth}\vspace*{3pt}}
\newcommand{\mpe}{\vspace*{3pt}\end{minipage}}

\newcommand{\resultul}[1]{\underline{\underline{#1}}}
\newcommand{\parskp}{$ $\\}	% new line after paragraph
\newcommand{\parnl}{\unskip\parfillskip 0pt \par} % justified new line (brace line!)

% mathematische Abkürzungen
\newcommand{\corr}{\;\widehat{=}\;}
\newcommand{\mdeg}{^{\circ}}
\newcommand{\nok}[2]{\binom{#1}{#2}}	% n über k BESSER: \binom{n}{k}
\newcommand{\mtr}[1]{\begin{pmatrix}#1\end{pmatrix}}	% Matrix
\newcommand{\dtr}[1]{\begin{vmatrix}#1\end{vmatrix}}	% Determinante (Betragsmatrix)
\LetLtxMacro{\originalVec}{\vec}
\renewcommand{\vec}[1]{\underline{#1}}	% Vektorschreibweise
\newcommand{\imptnt}[1]{\colorbox{red!30}{#1}}	% Wichtiges
\newcommand{\intd}[1]{\,\mathrm{d}#1}
\newcommand{\diffd}[1]{\mathrm{d}#1}
% für Module-Rechnung: \pmod{}
\newcommand{\unit}[1]{\,\mathrm{#1}}
%\LetLtxMacro{\ntilde}{\tilde}
%\renewcommand{\tilde}{\widetilde}

% weiter Abkürzungen
\newcommand{\gdw}{genau dann wenn}
\newcommand{\lecdate}[1]{\marginpar{\small\textcolor{gray}{Vorlesung\newline #1}}}

\faculty{Fakultät Informatik/Mathematik}
\chair{Lehrstuhl für Künstliche Intelligenz}
\subject{lecture}
\title{Neuroinformations\-verarbeitung}
\author{Falk-Jonatan Strube}
\professor{Prof. Dr. Hans-Joachim Böhme}

\slidePath{Vorlesung/}
\slideScale{.55}
\begin{document}
\maketitle
\tableofcontents

\chapter*{Einführung}
Klausur: Keine Unterlagen oder Hilfsmittel erlaubt

\setcounter{chapter}{-1}
\chapter{Biologischer Hintergrund}
\lecdate{10.10.2017}
\slides{NI_intro_bio}{1}
\section{Einordnung der Neuroinformatik}
Künstliche Intelligenz:
\begin{itemize}
\item Logik und Deduktion
\item Mustererkennung
\item Maschinelles Lernen
\item Spracherkennung
\item …
\item \emph{Neuroinformationsverarbeitung}
\begin{itemize}
\item Biologie
\item Psychologie
\item Computational Neuroscience
\item …
\item \emph{Neuronalen Netze}
\end{itemize}
\end{itemize}

\section{Geschichte der Neuroinformatik}
\slides{NI_intro_bio}{2}
\subsection{Lernen}
\begin{center}
\fbox{Biologisches System} $\longleftrightarrow$ \fbox{Umwelt}
\end{center}
\emph{Lernen}: Optimierung der „Fitness“\\
(Biologisches System soll besser in seiner Umwelt zurecht kommen).\\
Benötigt dafür: Funktion, die diese Fitness beschreibt.
\subsection{Klassifikation}
Einfache Klassifikation: Die, die links und die, die rechts von $\Theta$ liegen.
\begin{center}
\begin{tikzpicture}[scale=.5]
% Koordinatenachsen:
\draw [-latex] (-5.5,0) -- (7,0) node[below right]{$y$};
\draw [-latex] (0,-0.5) node[right]{$0$} -- (0,3.5) node[above left]{$z$};
% Funktion:
\draw [thick] (-3.5,0) -- (2.5,0) node[below]{$\Theta$} -- (2.5,2) -- (6,2);
\draw (0,2) -- (-0.2,2) node[left]{$1$};
\end{tikzpicture}
\end{center}
Eine Gerade als einfachste Form einer Diskriminanzfunktion (einer Punktwolke). Orthogonal zu dieser Geraden kann wieder die klassifizierende Funktion mit dem $\Theta$ an der Diskriminanzfunktion anlegen.
\begin{center}
\begin{tikzpicture}[scale=.5]
% Koordinatenachsen:
\draw [-latex] (-5.5,0) -- (7,0) node[below right]{$x_1$};
\draw [-latex] (0,-0.5) -- (0,3.5) node[above left]{$x_2$};
% Funktion:
\draw (-2,-1.5) -- (5.5,3.5);
% Punkte oben:
\draw [orange] (2,3) circle (0.2);
\draw [orange] (2.75,2.5) circle (0.2);
\draw [orange] (1,1.5) circle (0.2);
\draw [orange] (1.5,2) circle (0.2);
\draw [orange] (1,2.5) circle (0.2);
% Punkte unten
\draw [blue] (4.25,2) circle (0.2);
\draw [blue] (4.5,1) circle (0.2);
\draw [blue] (2.5,0.5) circle (0.2);
\draw [blue] (3.75,1) circle (0.2);
\draw [blue] (3.25,1.5) circle (0.2);
\end{tikzpicture}
\end{center}
Kontrast: XOR -- Welche lineare Funktion kann diese Klassen separieren? Keine!
\begin{center}
\begin{tikzpicture}[scale=1]
% Koordinatenachsen:
\draw [-latex] (-0.5,0) -- (2,0) node[below right]{$x_1$};
\draw [-latex] (0,-0.5) -- (0,2) node[above left]{$x_2$};
% Funktion:
% Punkte oben:
\draw [orange] (1,0) circle (0.1);
\draw [orange] (0,1) circle (0.1);
% Punkte unten
\draw [blue] (1,1) circle (0.1);
\draw [blue] (0,0) circle (0.1);
\end{tikzpicture}
\end{center}
Lösungsmöglichkeiten: Einteilung in Mengen -- die Diagonale und der Rest. 

\section{Das Gehirn}
\subsection{Anatomie}
\slides{NI_intro_bio}{3}
\subsection{Vorstellung der Funktionalität Anfang des 20. Jahrhunderts}
\slides{NI_intro_bio}{4}

\subsection{Neuronen}
\subsubsection{Experimenteller Nachweis von Einzelneuronen}
\slides{NI_intro_bio}{6}
\subsubsection{Strukturelle Vielfalt der Neuronentypen}
\slides{NI_intro_bio}{7}
\subsubsection{Strukturelle und funktionelle Details}
\slides[.7]{NI_intro_bio}{10}
Wichtige Bestandteile:
\begin{itemize}
\item Dendrit mit Kontaktstellen (Axon von anderer Zelle)\\
Kopplung durch Synapsen beschreibt man mit einem Gewichtsvektor $\vec{w}$, dessen Komponenten adaptierbar sind $\to$ lernen!
\item Axon\\
Jeder Eingang bildet \emph{eine} Komponente eines Eingabevektors $\vec{x}$.
\item Zellkörper\\
„verrechnet“ $\vec{x}$ und $\vec{w}$
\end{itemize}
Die Nervenzelle hat eine besondere Art der Zellmembran: Semipermeable Membran mit Ionen-Kanälen (Teilbild rechts oben).

\subsection{Medizinische Verfahren für funktionelle Untersuchungen}
\slides{NI_intro_bio}{8}

\subsection{Hirnaktivität bei Intellektuellen Aufgaben}
\slides{NI_intro_bio}{9}

\section{Neuronale Netze in der Robotik}
\subsection{Lokale Fahrzeugnavigation durch Experten-Cloning}
\slides{NI_intro_bio}{12}
\subsection[360 Grad Geräusch-Lokalisationssystems für Roboter]{Biologisch motiviertes Modell eines 360 Grad Geräusch-Lokalisationssystems für Roboter}
\slides{NI_intro_bio}{13}
\subsubsection*{Experimentelle Ergebnisse}
\slides{NI_intro_bio}{14}

\chapter{Grundlagen}
\slides{intro_bio}{1}
\section{Einordnung der Neuroinformatik}
Künstliche Intelligenz:
\begin{itemize}
\item Logik und Deduktion
\item Mustererkennung
\item Maschinelles Lernen
\item Spracherkennung
\item …
\item \emph{Neuroinformationsverarbeitung}
\begin{itemize}
\item Biologie
\item Psychologie
\item Computational Neuroscience
\item …
\item \emph{Neuronalen Netze}
\end{itemize}
\end{itemize}

\section{Geschichte der Neuroinformatik}
\slides{intro_bio}{2}
\subsection{Lernen}
\begin{center}
\fbox{Biologisches System} $\longleftrightarrow$ \fbox{Umwelt}
\end{center}
\emph{Lernen}: Optimierung der „Fitness“\\
(Biologisches System soll besser in seiner Umwelt zurecht kommen).\\
Benötigt dafür: Funktion, die diese Fitness beschreibt.
\subsection{Klassifikation}
Einfache Klassifikation: Die, die links und die, die rechts von $\Theta$ liegen.
\begin{center}
\begin{tikzpicture}[scale=.5]
% Koordinatenachsen:
\draw [-latex] (-5.5,0) -- (7,0) node[below right]{$y$};
\draw [-latex] (0,-0.5) node[right]{$0$} -- (0,3.5) node[above left]{$z$};
% Funktion:
\draw [thick] (-3.5,0) -- (2.5,0) node[below]{$\Theta$} -- (2.5,2) -- (6,2);
\draw (0,2) -- (-0.2,2) node[left]{$1$};
\end{tikzpicture}
\end{center}
Eine Gerade als einfachste Form einer Diskriminanzfunktion (einer Punktwolke). Orthogonal zu dieser Geraden kann wieder die klassifizierende Funktion mit dem $\Theta$ an der Diskriminanzfunktion anlegen.
\begin{center}
\begin{tikzpicture}[scale=.5]
% Koordinatenachsen:
\draw [-latex] (-5.5,0) -- (7,0) node[below right]{$x_1$};
\draw [-latex] (0,-0.5) -- (0,3.5) node[above left]{$x_2$};
% Funktion:
\draw (-2,-1.5) -- (5.5,3.5);
% Punkte oben:
\draw [orange] (2,3) circle (0.2);
\draw [orange] (2.75,2.5) circle (0.2);
\draw [orange] (1,1.5) circle (0.2);
\draw [orange] (1.5,2) circle (0.2);
\draw [orange] (1,2.5) circle (0.2);
% Punkte unten
\draw [blue] (4.25,2) circle (0.2);
\draw [blue] (4.5,1) circle (0.2);
\draw [blue] (2.5,0.5) circle (0.2);
\draw [blue] (3.75,1) circle (0.2);
\draw [blue] (3.25,1.5) circle (0.2);
\end{tikzpicture}
\end{center}
Kontrast: XOR -- Welche lineare Funktion kann diese Klassen separieren? Keine!
\begin{center}
\begin{tikzpicture}[scale=1]
% Koordinatenachsen:
\draw [-latex] (-0.5,0) -- (2,0) node[below right]{$x_1$};
\draw [-latex] (0,-0.5) -- (0,2) node[above left]{$x_2$};
% Funktion:
% Punkte oben:
\draw [orange] (1,0) circle (0.1);
\draw [orange] (0,1) circle (0.1);
% Punkte unten
\draw [blue] (1,1) circle (0.1);
\draw [blue] (0,0) circle (0.1);
\end{tikzpicture}
\end{center}
Lösungsmöglichkeiten: Einteilung in Mengen -- die Diagonale und der Rest. 

\section{Das Gehirn}
\subsection{Anatomie}
\slides{intro_bio}{3}
\subsection{Vorstellung der Funktionalität Anfang des 20. Jahrhunderts}
\slides{intro_bio}{4}

\subsection{Neuronen}
\subsubsection{Experimenteller Nachweis von Einzelneuronen}
\slides{intro_bio}{6}
\subsubsection{Strukturelle Vielfalt der Neuronentypen}
\slides{intro_bio}{7}
\subsubsection{Strukturelle und funktionelle Details}
\slides[.7]{intro_bio}{10}
Wichtige Bestandteile:
\begin{itemize}
\item Dendrit mit Kontaktstellen (Axon von anderer Zelle)\\
Kopplung durch Synapsen beschreibt man mit einem Gewichtsvektor $\vec{w}$, dessen Komponenten adaptierbar sind $\to$ lernen!
\item Axon\\
Jeder Eingang bildet \emph{eine} Komponente eines Eingabevektors $\vec{x}$.
\item Zellkörper\\
„verrechnet“ $\vec{x}$ und $\vec{w}$
\end{itemize}
Die Nervenzelle hat eine besondere Art der Zellmembran: Semipermeable Membran mit Ionen-Kanälen (Teilbild rechts oben).

\subsection{Medizinische Verfahren für funktionelle Untersuchungen}
\slides{intro_bio}{8}

\subsection{Hirnaktivität bei Intellektuellen Aufgaben}
\slides{intro_bio}{9}

\section{Neuronale Netze in der Robotik}
\subsection{Lokale Fahrzeugnavigation durch Experten-Cloning}
\slides{intro_bio}{12}
\subsection[360 Grad Geräusch-Lokalisationssystems für Roboter]{Biologisch motiviertes Modell eines 360 Grad Geräusch-Lokalisationssystems für Roboter}
\slides{intro_bio}{13}
\subsubsection*{Experimentelle Ergebnisse}
\slides{intro_bio}{14}

\chapter{Prinzipielle Anwendungsfälle von KNNs}
\section{Mathematische Notation}
\slides{WS2016_Kap1_Kap2}{4}
\subsection{Neuronale Bezeichner}
\slides{WS2016_Kap1_Kap2}{5}
$w_{ij}$: Sender $i$, Empfänger $j$
\slides{WS2016_Kap1_Kap2}{6}
$d$: Abstandsmaß (bspw. Euklidisch oder Hamming-Distanz)
\section{Neuronenmodell}
\slides{WS2016_Kap1_Kap2}{7}
\slides{WS2016_Kap1_Kap2}{8}
\unimptnt{
\section{Netzwerk- \& Lernparadigmen}
\slides{WS2016_Kap1_Kap2}{9}
}
\section{Lernparadigmen}
\subsection{Hebb'sche  Lernregeln}
Korrelationslernen, unüberwacht
\slides{WS2016_Kap1_Kap2}{10}
\begin{center}
\begin{tikzpicture}
\draw (0,0) node{$i$} circle (0.5) node[above = .5]{Sendeneuron};
\draw (4,0) node{$j$} circle (0.5) node[above = .5]{Empfangsneuron};
\draw [-hooks] (0.5,0) -- (2,0);
\draw [hooks-](2,0) -- (3.5,0);
\draw [decorate, decoration={brace, amplitude=5pt}](2.5,-0.5) -- (1.5,-0.5);
\node at (2,-1) {$w_{ij}$};
\end{tikzpicture}
\end{center}
\paragraph{Frage:} Wann soll sich der Gewichtswert der Verbindungssynapse $w_{ij}$ ändern?
\paragraph{Idee:} Gewicht proportional zur Aktivität prä- und postsynaptische Seite erhöhen.\bigskip\\
Als binäre Wertetabelle:\\
\begin{tabular}{c | c c}
$y_i / y_j$ & 0 & 1\\
\hline
0 & 0 & 0\\
1 & 0 & 1
\end{tabular}\\
$\Delta w_{ij} = \mu \cdot y_i \cdot y_j$ \qquad ($\mu$: Lernrate mit $0 < \mu \leq 1$)\\
$\to$ Gewicht der Synapse wird verstärkt, wenn \emph{beide} Neuronen aktiv sind.\\
$\Rightarrow$ Korrelations-LR\\
\emph{Nachteil:} keine Änderung $\Delta w_{ij}$ in negative Richtung möglich!

\unimptnt{
\subsection{SOFM}
Kohnen-Algorithmus, unüberwacht
\slides{WS2016_Kap1_Kap2}{11}
}

\subsection{Error-Driven}
überwacht
\slides{WS2016_Kap1_Kap2}{12}











\chapterN{Wiederholung Kap. 1+2}
\lecdate{24.10.2017}
\sectionN{Formales statisches Neuronenmodell}
$$\vec{x} = \begin{pmatrix}
x_1 \\ x_2 \\ \vdots \\ x_n
\end{pmatrix}
\To
\vec{w}_i=
\begin{pmatrix}
w_{1i}\\
w_{2i}\\
\vdots\\
w_{ni}
\end{pmatrix}
$$

$z_i=f(\vec{x}, \vec{w}_i$ … Zustand des Neurons $i$\\
$y_i=f(z_i)$ … Ausgabe des Neurons $i$

\subsectionN{Zustandsfunktionen}
\begin{anumerate}
\item auf Basis des Skalarprodukts\\
$z_i = w_{1i} \cdot x_1 + w_{2i}\cdot x_2 + \ldots + w_{ni}\cdot x_n - \Theta_i$\quad($\Theta_i$: Schwellwert)
\item auf Basis von Abstandsfunktionen\\
z.B. $z_i=\sqrt{(x_1-w_{1i})^2+\ldots+(x_n-w_{ni})^2}= L_2\text{-Norm bzw. euklidischer Abstand}$\\
Norm: $L_m=\sqrt[n]{(x_1-w_{1i})^m+\ldots+(x_n-w_{ni})^m}$\\
($L_1$-Norm wäre City-Block-Distanz)
\end{anumerate}

\subsectionN{Ausgabe- bzw. Transferfunktion}
\begin{enumerate}[label=(\arabic*)]
\item $y_i = f(z_i) = \frac{1}{1+e^{-z_i}}\To$ Fermi-  bzw. Signoidfunktion
\begin{center}
\begin{tikzpicture}
\begin{axis}[hide axis, axis equal, clip=false, scale=.5]
\addplot[domain=-3:3] {1/(1+e^(-x))}; 
\draw [-latex](-3.5,0) -- (3.5,0) node[below left]{$z_i$};
\draw [-latex](0,-0.1) -- (0,2) node[above right]{$y_i$}; 
\draw (0.1,1) -- (-0.1,1) node[left]{$1$};
\end{axis}
\end{tikzpicture}
\end{center}
\item $y_i = f(z_i)= \mathrm{sgn}(z_i) = \begin{cases}
0 & \text{für }z_i \leq 0\\
1 & \text{für }z_i >0
\end{cases}$
\item $y_i = f(z_i) = \exp \left( -\frac{z_i^2}{2 \sigma_i^2}\right)$
\begin{center}
\begin{tikzpicture}
\begin{axis}[hide axis, axis equal, clip=false, scale=.5]
\addplot[domain=0:3] {e^(-x^2/2)};
\draw [-latex](-0.5,0) -- (3.5,0) node[below left]{$z_i$};
\draw [-latex](0,-0.1) -- (0,2) node[above right]{$y_i$}; 
\draw (0.1,1) -- (-0.1,1) node[left]{$1$};
\draw [-latex] (0,0.5) -- (1,0.5) node [below left]{$\sigma_i$};
\end{axis}
\end{tikzpicture}
\end{center}
\end{enumerate}
Funktionen (1) und (2) setzen Skalarproduktaktivierfunktion voraus, Funktion (3) eine distanz-basierte Zustandsberechunng.

\subsectionN{Beispiel Skalarproduktaktivierfunktion}
Welche Funktion kann nur \underline{ein} Neuron mit Skalarproduktaktivierfunktion realisieren?
\begin{center}
\begin{tikzpicture}[scale=.7]
\draw  (0,0) circle (0.5) node{$2$};
\draw  (0,-2) circle (0.5) node{$1$};
\draw [-latex](-1.5,0) node[left]{$x_2$} -- (-0.5,0);
\draw [-latex](-1.5,-2) node[left]{$x_1$} -- (-0.5,-2);
\draw  (3,-1) circle (0.5) node{$i$};
\draw [-latex](0.5,0) -- node[above]{$w_{2i}$} (2.5,-0.5);
\draw [-latex](0.5,-2) -- node[below]{$w_{1i}$} (2.5,-1.5);
\draw [-latex](3.5,-1) -- (4.5,-1) node[right]{$y_i$};
\draw  (3,-3) circle (0.5) node{$B$};
\draw [-latex](3,-2.5) -- node[right]{$\Theta$} (3,-1.5);
\end{tikzpicture}\\
$B$: Bias-Neuron mit konstanter Ausgabe von $1$.
\end{center}
$$z_i=w_{1i}\cdot x_i + w_{2i}x_2-\Theta_i$$
$$y_i = f(z_i) = z_i$$
Beispieldaten (AND-Abbildung):
\begin{center}
\begin{tikzpicture}
\draw [-latex](-0.5,0) -- (1.5,0) node[below right]{$x_1$};
\draw [-latex](0,-0.1) -- (0,1.5) node[above left]{$x_2$}; 
\draw (0.1,1) -- (-0.1,1) node[left]{$1$};
\draw (1,0.1) -- (1,-0.1) node[below]{$1$};

\draw [orange] (0,1) circle (0.15);
\draw [orange] (0,0) circle (0.15);
\draw [orange] (1,0) circle (0.15);
\draw [blue] (1,1) circle (0.15);
\end{tikzpicture}
\end{center}
Neuron $i$ soll für $\vec{x}=(1,1)^T$ eine Ausgabe $>0$ und für $\vec{x}=(0,0)^T,\;(0,1)^T\;(1,0)^T$ Ausgabe $\geq 0$ liefern.\bigskip

Setze die Parameter $w_{1i}, \; w_{2i}$ und $\Theta_i$ so, dass die gewünschte Funktion entsteht:

z.B.: $w_{1i}=w_{2i}=1$; $\Theta=1$\\
Probe:\\
$\vec{x}^1=(0,0)^T$: $z_i=1\cdot 0 + 1 \cdot 0 -1 = -1\; \checkmark$\\
$\vec{x}^2=(0,1)^T$: $z_i=1\cdot 1 + 1 \cdot 0 -1 = \phantom{-}0\; \checkmark$\\
$\vec{x}^3=(1,0)^T$: $z_i=1\cdot 0 + 1 \cdot 1 -1 = \phantom{-}0\; \checkmark$\\
$\vec{x}^4=(1,1)^T$: $z_i=1\cdot 1 + 1 \cdot 1 -1 = \phantom{-}1\; \checkmark$\bigskip

Neuron bildet eine Lineare Diskriminanzfunktion aus ($\RR^2\To$Gerade, $\RR^2\To$Ebene, $\RR^n\To$ Hyperebene), die den Eingaberaum in zwei Hälften teilt.
\begin{center}
\begin{tikzpicture}
\draw [-latex](-0.5,0) -- (1.5,0) node[below right]{$x_1$};
\draw [-latex](0,-0.1) -- (0,1.5) node[above left]{$x_2$}; 
\draw (0.1,1) -- (-0.1,1) node[left]{$1$};
\draw (1,0.1) -- (1,-0.1) node[below]{$1$};

\draw [orange] (0,1) circle (0.15);
\draw [orange] (0,0) circle (0.15);
\draw [orange] (1,0) circle (0.15);
\draw [blue] (1,1) circle (0.15);
\draw [red](-0.4,1.4) -- (1.4,-0.4) node[below right]{Gerade für $z_i=0$};
\end{tikzpicture}
\end{center}
$\to$ Neuron realisiert eine Halbraumtrennung.\\
mit $y_i=\frac{1}{1+e^{-z_i}}$: 
\begin{center}
\begin{tikzpicture}[scale=.6]
\draw [-latex](-0.5,0) -- (1.5,0) node[below right]{$x_1$};
\draw [-latex](0,-0.1) -- (0,1.5) node[above left]{$x_2$}; 
\draw (0.1,1) -- (-0.1,1) node[left]{$1$};
\draw (1,0.1) -- (1,-0.1) node[below]{$1$};

\draw [orange] (0,1) circle (0.15);
\draw [orange] (0,0) circle (0.15);
\draw [orange] (1,0) circle (0.15);
\draw [blue] (1,1) circle (0.15);
\draw [red](-0.2,1.6) -- (1.6,-0.2);
\end{tikzpicture}
\end{center}

Offset/Schwelle ($\Theta$) gibt Abstand zum Ursprung an.

\subsectionN{Ermittlung optimale Parameter}
Wie ermittelt lineares Neuron sein optimalen Parameter?

Ausgangspunkt sei Hebb'sche Lernregel:
\begin{center}
\begin{tikzpicture}[scale=.7]
\draw (0,0) circle (0.5) node{$i$} node[above=.5]{Präsynapse}  node[below=.5]{Sender};
\draw (4,0) circle (0.5) node{$j$} node[above=.5]{Postsynapse}  node[below=.5]{Emfänger};
\draw [-latex] (0.5,0) -- node[above]{$w_{ij}$}(3.5,0);
\end{tikzpicture}
\end{center}
$$\Delta w_{ij} = \eta \cdot y_i \cdot y_j \quad \eta\text{: Lernrate}$$
Wir erweitern diese Lernregel und ersetzen den postsynaptischen Term $y_j$ durch einen \emph{Fehlerterm}!\\
Wir kennen zu jedem Eingabevektor die gewünschte Ausgabe (Zielausgabe bzw. \torange{Teachvorgabe}).\\
Am Bsp. der AND-Abbildung:\\
\begin{tabular}{c c | c}
$x_1$ & $x_2$ & \torange{$t$}\\\hline
0 & 0 & 0\\
0 & 1 & 0\\
1 & 0 & 0\\
1 & 1 & 1
\end{tabular}\\
Fehlerterm:\\
$(t-y_j)$ $\To$ Differenz zwischen tatsächlicher Ausgabe $y_j$ und gewünschter Ausgabe $t$.\\
Lernregel lautet jetzt:
$$\boxed{\Delta w_{ij}=\eta \cdot y_i \cdot (t-y_j)}$$
$\to$ \emph{Delta-Lernregel} bzw. \emph{Perceptron-Lernregel}\\
$(t-y_j)$: Fehler auf Ausgabeseite
\begin{center}
\begin{tikzpicture}[scale=.7]
\draw  (0,0) circle (0.5) node{$2$};
\draw  (0,-2) circle (0.5) node{$1$};
\draw [-latex](-1.5,0) node[left]{$x_2$} -- (-0.5,0);
\draw [-latex](-1.5,-2) node[left]{$x_1$} -- (-0.5,-2);
\draw  (3,-1) circle (0.5) node{$i$};
\draw [-latex](0.5,0) -- node[above]{$w_{2j}$} (2.5,-0.5);
\draw [-latex](0.5,-2) -- node[below]{$w_{1j}$} (2.5,-1.5);
\draw [-latex](3.5,-1) -- (4.5,-1) node[right]{$y_j$};
\draw  (3,-3) circle (0.5) node{$B$};
\draw [-latex](3,-2.5) -- node[right]{$\Theta_j$} (3,-1.5);
\draw [-latex](4.5,0.5) node[above]{$t$} -- (4.5,-0.5);
\draw [decorate, decoration={brace, amplitude=5pt}](5.5,1) -- node[right=.3]{$t-y_j$} (5.5,-1);
\draw [-latex] plot[smooth, tension=1] coordinates {(7,-0.5) (6,-2) (4,-2)};
\node at (6,-2)[below = .2, align=center]{Parameter-\\adaption};
\end{tikzpicture}
\end{center}
Mit dieser Lernregel durchläuft man den Trainingsdatensatz so lange, bis die Differenz zwischen $t$ und $y_i$ für alle Trainingsbeispiele minimal ist.

Diese Lernregel bildet die Basis für das Perceptron-Netzwerk (60er Jahre). Dies ist ein einschichtiges Netzwerk aus Sklarproduktneuronen: Beliebig viele $x_n$ werden von beliebig vielen $j_m$ verarbeitet.

Wie muss man sowohl das Netzwerk als auch die Lernregel erweitern, damit beliebig geformte Regionen voneinander separiert werden können?\\
$\To$ vom Perceptron zum Multi-Layer-Perceptron (MLP)\bigskip

\subsectionN{Beispiel XOR}
Bei Belegungen mit Anordnungen wie beim XOR:
\begin{center}
\begin{tikzpicture}[scale=.7]
\draw  (-2,0) circle (0.5) node{$1$};
\draw  (-2,-2) circle (0.5) node{$2$};
\draw [-latex](-3.5,0) node[left]{$x_1$} -- (-2.5,0);
\draw [-latex](-3.5,-2) node[left]{$x_2$} -- (-2.5,-2);
\draw  (0,0) circle (0.5) node{$3$};
\draw  (0,-2) circle (0.5) node{$4$};
\draw [-latex](-1.5,0)-- (-0.5,0);
\draw [-latex](-1.5,-2)-- (-0.5,-2);
\draw  (3,-1) circle (0.5) node{$5$};
\draw [-latex](0.5,0) -- (2.5,-0.5);
\draw [-latex](0.5,-2) -- (2.5,-1.5);
\draw [-latex](3.5,-1) -- (4.5,-1) node[right]{$y_5$};
\draw [-latex](-1.5,-2) -- (-0.5,-0.5);
\draw [-latex](-1.5,-0.5) -- (-0.5,-1.5);
\end{tikzpicture}
\end{center}
Neuronen 1, …, 5 sind Skalarprodukt-neuronen mit $y_i = \frac{1}{1+e{-z_i}}$ $\to$ stetig differenzierbare Funktion. Für unsere Darstellung der $y_4$-$y_5$-Ebene verwenden wir $y=f(z) = \begin{cases}
1 & \text{für } z> 0\\
0 & \text{für }z \leq 0
\end{cases}$

Wertetabelle:\\
\begin{tabular}{c c | c}
$x_1$ & $x_2$ & t\\\hline
0 & 0 & 0\\
0 & 8 & 1\\
8 & 0 & 1\\
8 & 8 & 0\\
\end{tabular}
\begin{center}
\begin{tikzpicture}[scale=1.5]
\draw [-latex, gray](-0.5,0) -- (1.5,0) node[below right]{$x_1$};
\draw [-latex, gray](0,-0.1) -- (0,1.5) node[above left]{$x_2$}; 
\draw [gray](0.1,1) -- (-0.1,1) node[left]{$8$};
\draw [gray](1,0.1) -- (1,-0.1) node[below]{$8$};

\draw [orange] (0,1) circle (0.15) node{$1$};
\draw [blue] (0,0) circle (0.15) node{$1$};
\draw [orange] (1,0) circle (0.15) node{$2$};
\draw [blue] (1,1) circle (0.15) node{$2$};
\draw (-0.5,0) -- (1,1.5) node[above right]{$G_3$};
\draw (0,-0.5) -- (1.5,1) node[above right]{$G_4$};
\end{tikzpicture}
\end{center}
Graphisches Ablesen/Ausrechnen der Parameter für Neuronen 3 und 4 liefert bspw.:\\
$w_{13}=1$; $w_{23}=-1$; $\Theta_3=-3$\\
$w_{14}=1$; $w_{24}=-1$; $\Theta_4=3$\\
$(0,0)^T$: $z_3=w_{13}x_1+w_{23}x_2-\Theta_3=3$; $y_3=1$, $z_4=w_{14}\ldots-\Theta_4=-3$; $y_4=0$\\
$(0,8)^T$: $z_3=-5 \to y_3=0$, $z_4=-11 \to y_4 = 0$\\
$(8,0)^T$: $z_4=12 \to y_3=1$, $z_4=5 \to y_4=1$\\
$(8,8)^T$: $z_3=3 \to y_3=1$, $z_4=-3 \to y_4 = 0$
\begin{center}
\begin{tikzpicture}[scale=1.5]
\draw [-latex, gray](-0.5,0) -- (1.5,0) node[below right]{$y_3$};
\draw [-latex, gray](0,-0.1) -- (0,1.5) node[above left]{$y_4$}; 
\draw [gray](0.1,1) -- (-0.1,1) node[left]{$8$};
\draw [gray](1,0.1) -- (1,-0.1) node[below]{$8$};
\draw [blue] (0,0) circle (0.15) node{$1$};
\draw [orange] (1,0) circle (0.2) node{$1/2$};
\draw [blue] (1,1) circle (0.15) node{$2$};
\draw (0,-0.5) -- (1.5,1) node[above right]{$G_5$};
\end{tikzpicture}
\end{center}









\chapter{Multilayer-Perceptron (MLP) und Error-Backpropagation-Algorithmus}
\lecdate{07.11.2017}
Aufbau:
\begin{center}
\begin{tikzpicture}[scale=.5]
\draw  (0,4) circle (0.5);
\draw  (0,1) circle (0.5);
\draw  (0,-0.5) circle (0.5);
\draw  (4,4) circle (0.5);
\draw  (4,1) circle (0.5);
\draw  (4,-0.5) circle (0.5);
\draw  (8,4) circle (0.5);
\draw  (8,1) circle (0.5);
\draw  (8,-0.5) circle (0.5);
\node at (0,4.5) [above, align=center] {Eingabe-\\schicht};
\node [above, align=center] at (4,4.5) {Hidden-\\schicht(en)};
\node [above, align=center] at (8,4.5) {Ausgabe-\\schicht};
\draw (0.5,4) -- (3.5,4);
\draw (4.5,4) -- (7.5,4);
\draw (0.5,-0.5) -- (3.5,-0.5);
\draw [thick, orange](4.5,-0.5) -- (7.5,-0.5);
\draw [thick, blue](0.5,3.5) -- (3.5,0);
\draw (0.5,0) -- (3.5,3.5);
\draw (4.5,3.5) -- (7.5,0);
\draw (4.5,0) -- (7.5,3.5);
\node at (4,2.5) {…};
\node at (0,2.5) {…};
\node at (8,2.5) {…};
\node at (0,-1) [below] {$k$};
\node at (4,-1) [below] {$j$};
\node at (8,-1) [below] {$i$};
\node at (-1.5,2) [left] {$\vec{x}^p \Rightarrow$};
\node at (11,4) [right] {$\vec{t}^p$};
\node at (10,2) {$\Rightarrow\vec{y}$};
\draw [-latex] (11,4) -- (10,3) -- (9,3);
\draw (11,3) -- (12.5,2) node[right, align=left]{Fehler kann für alle \\Ausgabeneuronen \\berechnet werden};
\end{tikzpicture}
\end{center}

\begin{itemize}
\item Feed-forward-network
\item vollvernetzt zwischen benachbarten Ebenen
\end{itemize}
Fehlerfunktion lautet:
$$\boxed{E^p = \frac{1}{2} \sum_{i=1}^{I}(t_i^p-y_i^p)^2}$$
(Quadratischer Ausgabefehler)\bigskip\\
Für alle Neuronen, ausgenommen die Eingabeneuronen, gilt:\\
$z_s = \sum w_{rs} \cdot y_s \To$ Skalarprodukt\\
$y_s = \frac{1}{1+e^{-z_s}}$\begin{center}
\begin{tikzpicture}
\begin{axis}[hide axis, axis equal, clip=false, scale=.7]
\addplot[domain=-2:2] {1/(1+e^(-4*x))}; 
\draw [-latex](-2.5,0) -- (2.5,0) node[below left]{$z_s$};
\draw [-latex](0,-0.1) -- (0,1.5) node[above right]{$y_s$}; 
\draw (0.1,1) -- (-0.1,1) node[left]{$1$};
\end{axis}
\end{tikzpicture}
\end{center}\bigskip
Für ausgezeichnetes Gewicht \torange{$w_{ji}$} würde mit Hilfe der Delta-Lernregel gelten:
$$\Delta w_{ji}=\eta \cdot y_j \cdot (t_i-y_i)$$
Wie adaptiert man das Gewicht \tblue{$w_{kj}$}?!

Idee: Die Delta-Lernregel mit Hilfe eines Gradientenverfahrens verallgemeinern, so dass sich ausgehend vom Fehler am NW-Ausgang \emph{alle} Parameter adaptieren lassen.\\
$\To$ wir notieren die Fehlerfunktion als Funktion, die vom zu adaptierenden Parameter abhängt (als Argument der Funktion). Zum Beispiel:
$$E^p(\vec{x}^p, \vec{W})=\torange{f(w_{ji})}$$
$$E^p(\vec{x}^p, \vec{W})=\tblue{f(w_{kj})}$$
$$E=\frac{1}{2}\sum(t_i-y_i)^2 =  \ldots =f(w_{ji})$$

Für alle Trainingsmuster/-beispiele soll der Fehler (am Ausgang) minimal sein. Dafür muss eben der Wert für die einzelnen Parameter der Übergänge gesucht werden, für die die Ausgabe einen minimalen Fehler hat.

Zum Bestimmen von $f(w_{ji})$ muss folgendes beachtet werden:

$\Delta w \approx -\frac{\partial E}{\partial w}$ (zum Suchen $w_{opt}$ in Richtung des Minimums bewegen)\\
$\to$ $\Delta w_{ji}\approx -\frac{\partial E}{\partial w-{ji}}$ \qquad $\Delta w_{kj}\approx - \frac{\partial E}{\partial w_{kj}}$\\
Wie hängt die Fehlerfunktion $E$ von Parameter $w_{ji}$ ab?\\
$\to$ $E_i=f(y_i)$ \quad $y_i=f(z_i)$ \quad $z_i=f(w_{ji})$\\
$\To E_i(y_i(z_i(w_{ji})))$\\
$\frac{\partial E_i}{\partial w_{ji}} = \frac{\partial E_i}{\partial y_i} \cdot \frac{\partial y_i}{\partial z_i} \cdot \frac{\partial z_i}{\partial w_{ji}}$ (Kettenregel)
\begin{itemize}
\item $\frac{\partial E_i}{\partial y_i}=\left[\frac{1}{2}(t_i-y_i)^2\right]_{y_i}=-(t_i -y_i)$
\item $\frac{\partial y_i}{\partial z_i}=\left[ \frac{1}{1+e^{-z_i}}\right]_{z_i}=\left[ (1+e^{-z_i})^{-1}\right]_{z_i}=-e^{-z_i}\cdot -(1+e^{-z_i})^{-2}\\
=\frac{e^{-z_i}}{(1+e^{-z_i})^2}=\frac{1}{1+e^{-z_i}}\cdot\frac{e^{-z_i}\torange{+1-1}}{1+e^{z_i}}=\frac{1}{1+e^{-z_i}}\cdot\left(\frac{1+e^{-z_i}}{1+e^{z_i}}-\frac{1}{1+e^{-z_i}}\right)=y_i(1-y_i)$
\item $\frac{\partial z_i}{\partial w_{ji}}=\left[ y_1^H \cdot w_{1i} + y_2^H \cdot w_{2i} + \ldots + y_j^H \cdot w_{ji} + \ldots \right]_{w_{ji}}=y_j^H$
\end{itemize}
Damit gilt:\\
$\frac{\partial E_i}{\partial w_{ji}}=-(t_i-y_i)\cdot y_i (1-y_i) \cdot y_i$\\
$\boxed{\Delta w_{ji}=\eta \left( - \frac{\partial E_i}{\partial w_{ji}} \right)}$\\
$\boxed{\Delta w_{ji} = \eta \cdot (t_i-y_i) \cdot y_i (1-y_i) \cdot y_j}$\\
Für $\Delta w_{kj}$ ergibt sich jetzt $E(y_i(z_i(y_j(z_j(w_{kj})))))$:\\
$\frac{\partial E}{\partial w_{jk}} = \frac{\partial E}{\partial y_i} \cdot \frac{\partial y_i}{\partial z_i} \cdot \frac{\partial z_i}{\partial y_j}\cdot \frac{\partial y_j}{\partial z_j}\cdot \frac{\partial z_j}{\partial w_{kj}}$\\
$\frac{\partial E}{\partial w_{jk}}= (-1) \sum_i (t_i-y_i) \cdot y_i(1-y_i) \cdot w_{ji} \cdot y_j(1-y_j) \cdot y_k$\\
Dabei sei $\boxed{ \delta_i =(t_i-y_i) \cdot y_i(1-y_i)}$ der Fehler am Knoten $i$.\\
Somit ist $\frac{\partial E}{\partial w_{jk}}=(-1) \sum_i \delta_i \cdot w_{ji} \cdot y_j(1-y_j) \cdot y_k$\\
Dabei sei $\sum_i \delta_i \cdot w_{ji} \cdot y_j(1-y_j) = \delta_j$ der Fehler am Knoten $j$.\\
Somit ist $\Delta w_{kj}=\eta \cdot y_j(1-y_j) \cdot y_k \sum_i \delta _i \cdot w_{ji} = \eta \cdot \delta_j \cdot y_k$\\
$\boxed{\Delta w_{ji}=\eta \cdot \delta_i \cdot y_j}$
$\boxed{\Delta w_{kj}=\eta \cdot \delta_j \cdot y_k}$ usw.






\section{Beispiel}
%%% Handout
Wichtig: An Neuron wird der $z$-Wert zusammengerechnet. Für die folgenden Neuronen ist dann der $y$-Wert relevant, der sich aus dem $z$-Wert ergibt.
\begin{align*}
z_5 &= 0,6\cdot 0 &&+ 0,3 \cdot 1 &&+ 0,1 \cdot 0 &&+ (-0,3) \cdot 1 &&+ 0 &&= 0\\
z_6 &= 0,1\cdot 0 &&+ 0,1 \cdot 1 &&+ 0,1 \cdot 0 &&+ 0,1 \cdot 1 &&-0,2 &&= 0\\
z_7 &= 0,4\cdot 0 &&+ (-0,1) \cdot 1 &&+ 0,1 \cdot 0 &&+ (-0,2) \cdot 1 &&+ 0,3 &&= 0
\end{align*}

\lecdate{14.11.2017}
Wichtig: Eingangsvariable $\vec{x}^p$ und Vorgabe $\vec{t}^p$.\\
Ergebnis im Vergleich zur Vorgabe:\\
$\vec{y}^p=\mtr{y_8\\y_9}=\mtr{0,5\\0,5} \Leftrightarrow \vec{t}^p = \mtr{t_8\\t_9}=\mtr{0\\1}$
\begin{center}
\begin{tikzpicture}
\draw  (-4,3) rectangle node[align=left]{Neuronales Netz\\\\ Parameter} (0,1);
\draw [-latex] (-5.5,2) node[left]{$\vec{x}$} -- (-4,2);
\draw [-latex] (0,2) -- (1.5,2) node[right]{$\vec{y}$};
\draw [-latex] (-1.5,4.5) node[left]{$E$} -- (-1,4.5) node[below right = .2, align=left]{Lernregel\\Trainingsalgorithmus} -- (0,0.5);
\end{tikzpicture}
\end{center}
$$\Delta w(t) = \eta \cdot \left( -\frac{\partial E}{\partial w}\right)$$
$$w(t+1) = w(t) + \Delta w(t)$$
Beispielweg von einem Eingang zum Ausgang:\\
$\Delta w_{58}=\eta \cdot \delta_8 \cdot y_5$ mit $\delta_8 = \underbrace{(t_8-y_8)}_{F'(y_8)\cdot (-1)} \cdot \underbrace{y_8 \cdot (1-y_8)}_{F'(z_8)}$\\
$\Delta w_{15}=\eta\cdot \delta_5 \cdot y_1$ mit $\delta_5 = (w_{58}\cdot \delta_8 + w_{59} \cdot \delta_9) \cdot y_5 \cdot (1-y_5)$\\
Damit lässt sich berechnen:
\subsection{Erster Trainingsschritt}
Gesucht: $\Delta \vec{w}_8$, $\Delta\vec{w}_9$, $\Delta\vec{w}_5$, $\Delta\vec{w}_6$, $\Delta\vec{w}_7$ (Training)
\begin{itemize}
\item $\Delta \vec{w}_8 = (\Delta w_{B8},\; \Delta w_{58},\; \Delta w_{68},\; \Delta w_{78})^T$
\begin{itemize}
\item $\Delta w_{B8} = 0,1 \cdot (0-0,5) \cdot 0,5 \cdot (1-0,5) \cdot 1 = -0,0125$
\item $\Delta w_{58} = 0,1 \cdot \underbrace{(0-0,5)\cdot 0,5\cdot (1-0,5)}_{\delta_8=-0,125} \cdot 0,5 = -0,00625$
\item $\Delta w_{68} = \Delta w_{58}$
\item $\Delta w_{78} = \Delta w_{58}$
\end{itemize}
\item $\Delta \vec{w}_9 = (\Delta w_{B9},\; \Delta w_{59},\; \Delta w_{69},\; \Delta w_{79})^T$
\begin{itemize}
\item $\Delta w_{B9} = 0,0125$
\item $\Delta w_{59} = 0,1 \cdot \underbrace{(1-0,5) \cdot 0,5 \cdot (1-0,5)}_{\delta_9=0,125} \cdot 0,5 = 0,00625$
\item $\Delta w_{69} = \Delta w_{59}$
\item $\Delta w_{79} = \Delta w_{59}$
\end{itemize}
\item $\Delta \vec{w}_5 = (w_{B5},\; w_{15},\; w_{25},\; w_{35},\; w_{45})^T$
\begin{itemize}
\item $\Delta w_{B5} = 0,1 \cdot \underbrace{[(-0,6) \cdot (-0,125) + (-0,6) \cdot 0,125]}_{=0}\cdot 0,5 \cdot (1-0,5) \cdot 0 = 0$
\item $\Delta w_{15} = 0$
\item $\Delta w_{25} = 0$
\item $\Delta w_{35} = 0$
\item $\Delta w_{45} = 0$
\end{itemize}
\item $\Delta \vec{w}_6 = (w_{B6},\; w_{16},\; w_{26},\; w_{36},\; w_{46})^T$
\begin{itemize}
\item $\Delta w_{B6} = 0,1 \cdot [0,2 \cdot (-0,125) + (-0,4) \cdot 0,125]\cdot 0,5 \cdot (1-0,5) \cdot 1 = -0,001875$
\item $\Delta w_{16} = 0$
\item $\Delta w_{26} = \Delta w_{B6}$
\item $\Delta w_{36} = 0$
\item $\Delta w_{46} = \Delta w_{B6}$
\end{itemize}
\item $\Delta \vec{w}_7 = (w_{B7},\; w_{17},\; w_{27},\; w_{37},\; w_{47})^T$
\begin{itemize}
\item $\Delta w_{B7} = 0,1 \cdot [(-0,6) \cdot (-0,125) + (-0,4) \cdot 0,125]\cdot 0,5 \cdot (1-0,5) \cdot 1 = -0,000625$
\item $\Delta w_{17} = 0$
\item $\Delta w_{27} = \Delta w_{B7}$
\item $\Delta w_{37} = 0$
\item $\Delta w_{47} = \Delta w_{B7}$
\end{itemize}
\end{itemize}
Optionen beim Korrigieren: Entweder Änderungen direkt drauf rechnen (direct learning), oder als zusätzlichen „Delta-Bias“ in einer Variable speichern und dann drauf rechnen (batch learning).

Warum initialisieren wir die Parameter mit \emph{zufälligen} Werten?
% Folien: Übung, Folien 5+6
\begin{itemize}
\item Damit die Änderungen möglichst nicht klein werden.
\end{itemize}
Was passiert, wenn alle Parameter mit Null initialisiert werden?
\begin{itemize}
\item $z_5=z_6=z_7=0$
\item $y_5=y_6=y_7=0.5$
\item $z_8=z_9=0$
\item $y_8=y_9=0.5$
\end{itemize}
$\to$ die $y$-Werte bleiben gleich.\\
Beim Training:\\
1.Schritt:
\begin{itemize}
\item $\Delta \vec{w}_8$ und $\Delta \vec{w}_9$ bleiben wie bei der zufälligen Verteilung.
\item $\Delta \vec{w}_5=\Delta \vec{w}_6=\vec{w}_7$ ist dann $\vec{0}$.
\end{itemize}
Damit:
\begin{itemize}
\item $\vec{w}_8(t+1)=\vec{w}_8(t)+\Delta w_8(t) = (-0,0125, -0,00625,\ldots)^T$
\item $\vec{w}_9(t+1)=\vec{w}_9(t)+\Delta w_9(t) = (0,0125, 0,00625,\ldots)^T$
\item $\vec{w}_4(t+1) = \vec{w}_4(t)$, gleiches für $\vec{w}_5$ und $\vec{w}_6$
\end{itemize}
Und somit:
\begin{itemize}
\item $z_5=z_6=z_7= 0$
\item $y_5=y_6=y_7 = 0.5$
\item $z_8=-0.0125\cdot 1 + 3 \cdot ((-0.00625)\cdot  0.5)=-0.021875$
\item $z_9=0.0125\cdot 1 + 3 \cdot (0.00625\cdot  0.5)=0.021875$
\item $y_8=0.5055$
\item $y_9=0.4945$
\end{itemize}
2. Schritt:
\begin{itemize}
\item $\Delta\vec{w}_8=(-0.0124,-0.0062,-0.0062,\ldots)^T$
\item $\Delta\vec{w}_9=(0.0124, 0.0062,0.0062,\ldots)^T$
\item $\Delta \vec{w}_5=\Delta \vec{w}_6 = \Delta\vec{w}_7 = (0.00004, 0, 0.00004, 0, 0.00004)^T$
\end{itemize}
Damit:
\begin{itemize}
\item $\vec{w}_5(t+2)=\vec{w}_6(t+2)=\vec{w}_7(t+2)=(0.00004, 0, 0.00004, 0, 0.00004)$
\end{itemize}
$\to$ alle Hidden-Knoten sind gleich $\to$ man könnte die Hidden-Knoten auch auf einen reduzieren. Damit kann aber bspw. nur eine Gerade abgebildet werden (Vergleich XOR -- braucht 2 Geraden) (vgl. Symmetriebruch).\\
$\To$ man braucht unterschiedliche/zufällige Initialwerte, damit \emph{unterschiedliche} Geraden angelernt werden können!

\section{Multi Layered Perceptron}
\slides{NI_WS2017_Kap3_MLP}{1}

\section{Grundprinzip Training}
\slides{NI_WS2017_Kap3_MLP}{2}

\section{Error-Backtracking-Algorithmus}
\slides{NI_WS2017_Kap3_MLP}{3}
\subsection{Gradientanstieg}
\slides{NI_WS2017_Kap3_MLP}{4}
\subsection{Rückpropagierung des Fehlers am Netzwerk-Ausgang}
\slides{NI_WS2017_Kap3_MLP}{5}
\slides{NI_WS2017_Kap3_MLP}{6}
\subsection{Adaptionsvorschriften}
\slides{NI_WS2017_Kap3_MLP}{7}
\subsection{Fazit}
\slides{NI_WS2017_Kap3_MLP}{8}



\chapter{Erweiterung Error-Backpropagation-Algorithmus}
\lecdate{21.11.2017}
\section{Probleme des EBP-Algorithmus}
\subsection{Resultierend aus Gradientenabstieg}
\slides{NI_WS2016_Kap4_BP_ext}{4}
Kommt aus lokalem Minimum nicht mehr raus, weil dort der Gradient $0$ ist $\Rightarrow$ es kann keine Änderung mehr statt finden. Einzige Lösung: neu Initialisieren bzw. willkürlicher Schritt.
\slides{NI_WS2016_Kap4_BP_ext}{5}
\slides{NI_WS2016_Kap4_BP_ext}{6}
\slides{NI_WS2016_Kap4_BP_ext}{7}

\subsection{Resultierend aus Netzwerk-Topologie}
\slides{NI_WS2016_Kap4_BP_ext}{8}
Sehr weit außerhalb des Arbeitsbereichs bekommt man einen kleinen Gradienten ($\corr$ Adaptionsgeschwindigkeit gering)
\slides{NI_WS2016_Kap4_BP_ext}{9}

\section{Training mit Momentum-Term}
\slides{NI_WS2016_Kap4_BP_ext}{10}
\slides{NI_WS2016_Kap4_BP_ext}{11}
„Langsamer ins Tal gleiten, Schwung von starkem Abstieg ins Plateau mitnehmen.“

\section{Modifikation der Ableitung der Transferfunktion}
\slides{NI_WS2016_Kap4_BP_ext}{12}
Schrittgröße hat ein festes Minimum, Gradient gibt dann nur die Richtung an.

\section{Training mit Weight-Decay}
\slides{NI_WS2016_Kap4_BP_ext}{13}
\slides{NI_WS2016_Kap4_BP_ext}{14}
\slides{NI_WS2016_Kap4_BP_ext}{15}
\slides{NI_WS2016_Kap4_BP_ext}{16}

\section{Manhatten-Training}
\slides{NI_WS2016_Kap4_BP_ext}{17}
\slides{NI_WS2016_Kap4_BP_ext}{18}
$\To$ wie Modifikation der Ableitung der Transferfunktion, nur für beide Richtungen.

\section{Eigene Schrittweite für jeden Parameter}
\slides{NI_WS2016_Kap4_BP_ext}{19}
\slides{NI_WS2016_Kap4_BP_ext}{20}

\section{Resilient Backpropagation (Rprop)}
\slides{NI_WS2016_Kap4_BP_ext}{24}
\slides{NI_WS2016_Kap4_BP_ext}{25}
\slides{NI_WS2016_Kap4_BP_ext}{26}
\slides{NI_WS2016_Kap4_BP_ext}{27}
\slides{NI_WS2016_Kap4_BP_ext}{28}

\section{Quickpropagation (Quickprop)}
\slides{NI_WS2016_Kap4_BP_ext}{21}
\slides{NI_WS2016_Kap4_BP_ext}{22}
\slides{NI_WS2016_Kap4_BP_ext}{23}

\section{Fazit}
\slides{NI_WS2016_Kap4_BP_ext}{29}
\slides{NI_WS2016_Kap4_BP_ext}{30}







\chapter{Radial Basis Function-Netzwerk (RBF-Netzwerk)}
\section{Motivation}
\slides{NI_WS2016_Kap5_RBF}{1}
Datenpunkte als Prototypen der Abbildung $\to$ diese Umkreisen.
\begin{center}
\begin{tikzpicture}[scale=1]
% Koordinatenachsen:
\draw [-latex] (-0.5,0) -- (2,0) node[below right]{$x_1$};
\draw [-latex] (0,-0.5) -- (0,2) node[above left]{$x_2$};
% Funktion:
% Punkte oben:
\draw [orange] (1,0) circle (0.1);
\draw [orange] (0,1) circle (0.1);
\draw [red] (1,0) circle (0.3);
\draw [red] (0,1) circle (0.3);
% Punkte unten
\draw [blue] (1,1) circle (0.1);
\draw [blue] (0,0) circle (0.1);
\draw [red] (1,1) circle (0.3);
\draw [red] (0,0) circle (0.3);
\end{tikzpicture}
\end{center}
Kombiniere die Ausgaben des Netzwerks aus den Ausgaben der Prototyp- bzw. Referenz-Neuronen.\\
Jedes Neuron der Hidden-Schicht bildet einen (Radius von einem) Datenpunkt dar:
\begin{center}
\begin{tikzpicture}[scale=.5]
\draw  (0,1) circle (0.5);
\draw  (0,-1) circle (0.5);

\draw [red] (4,3) circle (0.5);
\draw [red] (4,1) circle (0.5);
\draw [red] (4,-1) circle (0.5);
\draw [red] (4,-3) circle (0.5);

\draw  (8,1) circle (0.5);
\draw  (8,-1) circle (0.5);
\draw (0.5,1) -- (3.5,3);
\draw (0.5,1) -- (3.5,-3);
\draw (0.5,-1) -- (3.5,3);
\draw (0.5,-1) -- (3.5,-3);
\draw (4.5,3) -- (7.5,1) -- (4.5,-3);
\draw (4.5,3) -- (7.5,-1) -- (4.5,-3);
\node at (3,0) {$\vdots$};
\node at (5,0) {$\vdots$};
\end{tikzpicture}
\end{center}
Skalarprodukt:
$$z=\vec{w}^T\cdot \vec{w}$$
Distanzbasierte Aktionsgerade:
$$z=f(\vec{x},\vec{w})$$
$\to$ Vektornormen\\
$\to$ $L_2$-Norm \\
$\to$ $\boxed{z=\sqrt{(x_1-w_1)^2+\ldots + (x_n-w_n)^2}}$ (Euklidischer Abstand) … $z$ als Abstand zwischen $\vec{x}$ und $\vec{w}$.
$$y=f(z)\approx e^{-z}$$
$$\boxed{f=\exp \left( -\frac{z^2}{2 \sigma^2} \right)} \text{ Gaußfunktion}$$
\slides{NI_WS2016_Kap5_RBF}{2}

\section{Topologie}
\slides{NI_WS2016_Kap5_RBF}{3}
\section{Berechnungsvorschriften}
\slides{NI_WS2016_Kap5_RBF}{4}
\slides{NI_WS2016_Kap5_RBF}{5}
\label{berechnung_rbf}
\section{Trainingsregime}
\slides{NI_WS2016_Kap5_RBF}{6}

\subsection{Bestimmung der Gewichte von der Eingangs- zur RBF-Schicht}
\slides{NI_WS2016_Kap5_RBF}{7}
\slides{NI_WS2016_Kap5_RBF}{8}

\subsection{Bestimmung der Varianzen}
\slides{NI_WS2016_Kap5_RBF}{9}
\slides{NI_WS2016_Kap5_RBF}{10}

\subsection{Bestimmung der Gewichte von der RBF- zur Ausgabeschicht}
\slides{NI_WS2016_Kap5_RBF}{11}
$\vec{A}=f(\vec{X}, \vec{W}, \vec{\sigma})$
\slides{NI_WS2016_Kap5_RBF}{12}
Achtung bei Inverser: Elemente dürfen nicht linear abhängig sein!

\subsection{Iteratives Nachtraining}
\slides{NI_WS2016_Kap5_RBF}{13}

\subsection{Weitere mathematische Betrachtungen}
\slides{NI_WS2016_Kap5_RBF}{14}

\section{Notizen}
\lecdate{05.12.2017}
\begin{itemize}
\item Genauigkeit des Netzwerks ist geringer, je weiter die Trainingspunkte auseinander liegen (je kleiner damit die Einzugsbereiche werden).
\item Gelerntes Netzwerk bildet nur Daten in der Nähe der gelernten Punkte ab. Das kann bspw. gut sein, wenn die Frage ist, wann der Arbeitsbereich verlassen wird.
\item Bei RBF gibt es immer nur eine Hidden-Schicht!
\end{itemize}

\section{Berechnungsbeispiel}
\subsection{Gegeben}
Abbildung von 4 Merkmalen/Eingängen auf 2 Klassen.
\begin{center}
\begin{tikzpicture}[scale=.5]
\draw  (0,0) circle (0.5) node[below=.5]{in};
\draw  (0,2) circle (0.5);
\draw  (0,4) circle (0.5);
\draw  (0,6) circle (0.5);

\draw  (8,2) circle (0.5) node[below=.5]{out};
\draw  (8,4) circle (0.5);

\node at (4,3) [align=center]{RBF-\\Schicht};
\draw [-latex] (1,3) -- (2.5,3);
\draw [-latex] (5.5,3) -- (7,3);
\end{tikzpicture}
\end{center}
$$\vec{x}^1=\mtr{0\\1\\0\\1} \quad \vec{x}^2=\mtr{0\\0\\1\\0} \quad \vec{x}^3=\mtr{1\\1\\0\\0}\quad \vec{t}^1=\mtr{0\\1} \quad \vec{t}^2 =\mtr{1\\0} \quad \vec{t}^3=\mtr{0\\1}$$
\subsection{Bestimmung der RBF-Zentren}
Einfacher Ansatz: für jeden Trainingspunkt ein Zentrum.\\
Besser: Ansatz mit möglichst geringem Quantisierungsfehler.
\subsubsection{Quantisierungsfehler}
Quantisierung: Einteilung\\
$\To$ Quantisierung eines Vektorraumes durch $m$ Referenz- bzw. Codebook-Vektoren.\\
$\To$ Bestimme $\forall \vec{x} \in X$ die Differenz zwischen dem am besten passenden Referenzvektor und $\vec{x}$.\\
Referenzvektoren mit Varianz $\sigma^2 \forall$ Referenzvektoren identische Varianz angenommen:\\
Jeder Punkt hat Gaußsche Verteilung um sich. Damit ist das am besten passende (Best-Matching-Neuron) das, welches am nahsten am Punkt dran liegt. Die Differenz ist somit die Distanz (i.d.R. euklidischer Abstand) zwischen dem am besten passenden Neuron und dem betrachteten Punkt. \\
Mit $BM=(w_1,w_2)^T$ und betrachteter Punkt $\vec{x}=(x_1,x_2)^T$ ist die Differenz:
$$d(\vec{x}, BM)=\sqrt{(x_1-w_1)^2+(x_2-w_2)^2}$$
Der Einzugsbereich eines Neuron ist durch den halben Abstand zu seinen Nachbarneuronen begrenzt (Delanney-Triangulation). Den daraus erhaltenen Einzugsbereich nennt man auch Voronoi-Region.

\subsubsection{Bestimmung}
\begin{itemize}
\item Definiere einen maximalen Quantisierungsfehler $d_{max}^2=1$, der nicht überschritten werden darf.\\
$\to d^2(\vec{x}, \vec{w}^{I,RBF}) \leq 1 \;\forall x$
\item Ordne $\vec{w}_5^{I,RBF}$ Vektor $\vec{x}^1$ zu:
$$\vec{w}_5^{I,RBF}=(0,1,0,1)^T$$
\item Prüfe, ob $\vec{x}^2$ und $\vec{x}^3$ durch Referenzvektor $\vec{w}_5^{I,RBF}$ abgedeckt sind:\\
$d(\vec{x}^2,\vec{w}_5^{I,RBF}) = 3 \quad > d_{max} \;\checkmark$\\
$d(\vec{x}^3,\vec{w}_5^{I,RBF}) = 2 \quad > d_{max} \;\checkmark$
\item Ordne $\vec{w}_6^{I,RBF}$ Vektor $\vec{x}^2$ zu:
$$\vec{w}_6^{I,RBF}=(0,0,1,0)^T$$
\item Prüfe, ob $\vec{x}^3$ durch $\vec{w}_6^{I,RBF}$ abgedeckt ist:\\
$d(\vec{x}^3,\vec{w}_6^{I,RBF}) = 3 \quad > d_{max} \;\checkmark$
\item Ordne $\vec{w}_7^{I,RBF}$ Vektor $\vec{x}^3$ zu:
$$\vec{w}_7^{I,RBF}=(1,1,0,0)^T$$ 
\end{itemize}
Damit: 
\begin{center}
\begin{tikzpicture}[scale=.5]
\draw  (0,6) node{1} circle (0.5);
\draw  (0,4) node{2} circle (0.5);
\draw  (0,2) node{3} circle (0.5);
\draw  (0,0) node{4} circle (0.5) node[below=.5]{in};

\draw  (8,4) node{8} circle (0.5);
\draw  (8,2) node{9} circle (0.5) node[below=.5]{out};

\draw  (4,5) node{5} circle (0.5);
\draw  (4,3) node{6} circle (0.5);
\draw  (4,1) node{7} circle (0.5);

\draw (0.5,6) -- (3.5,5);
\draw (0.5,4) -- (3.5,5);
\draw (0.5,2) -- (3.5,5);
\draw (0.5,0) -- (3.5,5);

\draw (0.5,6) -- (3.5,3);
\draw (0.5,4) -- (3.5,3);
\draw (0.5,2) -- (3.5,3);
\draw (0.5,0) -- (3.5,3);

\draw (0.5,6) -- (3.5,1);
\draw (0.5,4) -- (3.5,1);
\draw (0.5,2) -- (3.5,1);
\draw (0.5,0) -- (3.5,1);
\draw [rotate=20] (4.5,3.5) ellipse (0.2 and .75) node[above=.4]{$\vec{x}^1$};

\draw (4.5,5) -- (7.5,4);
\draw (4.5,3) -- (7.5,4);
\draw (4.5,1) -- (7.5,4);

\draw (4.5,5) -- (7.5,2);
\draw (4.5,3) -- (7.5,2);
\draw (4.5,1) -- (7.5,2);

\draw (8.5,4) -- (9.5,4) node[right]{$y_8$};
\draw (8.5,2) -- (9.5,2) node[right]{$y_9$};
\draw [dashed] (5.5,7.5) circle (0.5);
\draw [dashed] (5.5,7) -- (6,5);
\draw [dashed](5.5,7) -- (6.5,5.5);
\end{tikzpicture}
\end{center}

\subsection{Ausgabevektor}
Weiterhin gegeben:
$$\vec{w}_8^{RBF,o}=\mtr{0\\1\\1} \quad \vec{w}_9^{RBF,o}=\mtr{0\\1\\0} \quad \sigma^2 = 0.5$$
Gesucht: $\vec{y}^{o}=\mtr{y_8\\y_9}$ für Eingabevektor $\vec{x}^1$\\
Rechnung: (siehe \autoref{berechnung_rbf})
\begin{itemize}
\item $z_5 = \sqrt{(x_1^1-w_{15})^2+\ldots+(x_4^1-w_{45})^2}=0$
\item $y_5=\exp\left(- \frac{z_5^2}{2 \sigma^2}\right)=1$
\item $z_6= \sqrt{(x_1^1-w_{16})^2+\ldots+(x_4^1-w_{46})^2}=\sqrt{3}$
\item $y_6 = e^{-3}$
\item $z_7 = \sqrt{(x_1^1-w_{17})^2+\ldots+(x_4^1-w_{47})^2}=\sqrt{2}$
\item $y_7 = e^{-2}$
\item $y_{RBF}=\left(1, e^{-3}, e^{-2}\right)^T$
\end{itemize}
Damit:
\begin{itemize}
\item $y_8=z_8=\vec{w}_8^T \cdot \vec{y}_{RBF} = e^{-3}+e^{-2}$
\item $y_9=z_9=\vec{w}_9^T \cdot \vec{y}_{RBF} = e^{-3}$
\end{itemize}
\subsection{Lernschritt}
Mittels Delta-Lernregel ($\eta=0,3$) einen Lernschritt für das Gewicht $w_{58}$ ausführen. Gesucht: $\Delta w_{58}$










\chapter{Learning Vector Quantization (LVQ)}
\section{Motivation}
\lecdate{09.01.2018}
\paragraph{Ziel}
Repräsentation der Eingabedaten erzeugen, die
\begin{itemize}
\item sowohl deren Statistik
\item als auch deren Klassenzugehörigkeit
\end{itemize}
berücksichtigt.

\paragraph{Vorgehen} Anzahl $M$ Referenzvektoren sollen eine Anzahl $P$ Datenpunkte repräsentieren/abbilden:
\begin{center}
\begin{tikzpicture}[scale=.5]
\draw [-latex] (-.5,0) -- (8,0) node[right]{$x_1$};
\draw [-latex] (0,-.5) -- (0,8) node[right]{$x_2$};
\draw [blue] plot[smooth cycle, tension=.7] coordinates {(1.5,6) (1,3.5) (3,4) (4,5) (4,6.5)};
\draw [orange] plot[smooth cycle, tension=.7] coordinates {(6.5,6.5) (5,5.5) (5,3) (8,3) (8,5.5)};
\draw [green] plot[smooth cycle, tension=.7] coordinates {(0.5,1.5) (3.5,3) (6,1) (3,0.5)};
\end{tikzpicture}
(3 Klassen)
\end{center}
\begin{itemize}
\item Definiere eine feste Anzahl von Referenzvektoren pro Klasse.
\item Initialisiere diese mit zufällig ausgewählten Beispielen der jeweiligen Klassen. $\to$ bspw. 4 pro Klasse:
\end{itemize}
\begin{center}
\begin{tikzpicture}[scale=.5]
\draw [-latex] (-.5,0) -- (8,0) node[right]{$x_1$};
\draw [-latex] (0,-.5) -- (0,8) node[right]{$x_2$};
\draw [blue] plot[smooth cycle, tension=.7] coordinates {(1.5,6) (1,3.5) (3,4) (4,5) (4,6.5)};
\draw [orange] plot[smooth cycle, tension=.7] coordinates {(6.5,6.5) (5,5.5) (5,3) (8,3) (8,5.5)};
\draw [green] plot[smooth cycle, tension=.7] coordinates {(0.5,1.5) (3.5,3) (6,1) (3,0.5)};
\draw [blue] (2,6) circle (0.2);
\draw [blue] (1.5,5.5) circle (0.2);
\draw [blue] (3.5,5) circle (0.2);
\draw [blue] (1.5,4) circle (0.2);
\draw [orange] (6.5,6) circle (0.2);
\draw [orange] (7,5.5) circle (0.2);
\draw [orange] (5.5,3.5) circle (0.2);
\draw [orange] (7.5,3.5) circle (0.2);
\draw [green] (1,1.5) circle (0.2);
\draw [green] (5.5,1) circle (0.2);
\draw [green] (4.5,1) circle (0.2);
\draw [green] (5,1.5) circle (0.2);
\end{tikzpicture}
\end{center}
\paragraph{Lernverfahren} 
\begin{itemize}
\item Versuche, die Lage der Referenzvektoren zu optimieren.
\item Nach Optimierung ist die Klassifikation von unbekannten Eingangsdaten möglich:\\
Das Klassenattribut des \emph{Best-Matching-Neurons} (-Referenzvektors) wird dem Eingangsdatum $\vec{x}$ zugeordnet (Zuweisung zum Referenzvektor, der am nächsten zum Eingangsvektor liegt).
\end{itemize}

\section{Optimierverfahren}
Für LVQ existieren \emph{viele} unterschiedliche Varianten.

\subsection{LVQ 1}
In jedem Adaptionsschritt  wird \emph{nur} das Best-Matching-Neuron (im Bezug zu der Menge der Beispiel- bzw.Trainingsdaten) angepasst.\\
Format der Trainingsdaten: $\vec{x}^{n}$ Label $k$\quad($k \in $ mögliche Klassen [in Beispiel $\tblue{1}, \torange{2}, \tgreen{3}$])

\subsubsection*{Lernalgorithmus}
\begin{enumerate}
\item Wähle ein $\vec{x}$ (aus Trainingsdaten).
\item Bestimme Best-Matching-Neuron $b$ nach Lernvorschrift:
$$\vec{w}_b=\underset{i=1\ldots P}{\operatorname{argmin}} \Vert \vec{x}-\vec{w}_i \Vert$$
Sowohl $\vec{x}$ als auch $\vec{w}_b$ verfügen über ihr jeweiliges Klassenlabel/-attribut.
\item Zwei Fälle abhängig davon, ob die Klassenattribute gleich sind oder nicht:
\begin{itemize}
\item Falls $\vec{x}$ und $\vec{w}_b$ gleiches Klassenattribut aufweisen:
$$\vec{w}_b(t+1) = \vec{w}_b(t)\tred{+}\eta(t)\left[ \vec{x}(t)-\vec{w}_b(t)\right]$$
$\To$ Verschiebung von $\vec{w}_b$ in Richtung $\vec{x}$.
\item Falls $\vec{x}$ und $\vec{w}_b$ verschiedene Klassenattribute aufweisen: 
$$\vec{w}_b(t+1) = \vec{w}_b(t)\tred{-}\eta(t)\left[ \vec{x}(t)-\vec{w}_b(t)\right]$$
$\To$ Verschiebung von $\vec{w}_b$ von $\vec{x}$ weg.
\end{itemize}
Zeitabhängige Lernrate $\eta(t)$ ist mit $\eta(t+1)=\eta(t)\cdot \alpha$ und $\eta(t=0)\approx 0.5$ definiert. Dabei liegt $0<\alpha\leq 1$, typische Größen sind $0.9,\; 0.99$ usw.\\
Durch die zeitabhängige Lernrate „friert“ die Verteilung irgendwann ein (und stellt somit den optimalen Zustand dar).\\
Wichtig: $t$ der Lernrate meint die gesamte Trainingsepoche (ein Mal mit allen Trainingsdaten trainiert).\\
Anmerkung: $\alpha=1$ wäre nur sinnvoll, wenn die Lernrate $\eta(t=0)$ sehr klein ist.
\end{enumerate}

\subsection{LVQ 2.1}
Im Unterschied zu LVQ 1:
\begin{itemize}
\item Nutzt sowohl Best-Matching-Neuron $b$ a-s auch Second-Best-Matching-Neuron $s$.
\item $\vec{x}$ muss in einem „Fenster“ bzw. zwischen $\vec{w}_b$ und $\vec{w}_s$ liegen:
\begin{center}
\begin{tikzpicture}[scale=.5]
\draw  plot[smooth, tension=.7] coordinates {(-2.5,4.5) (-2.5,2) (-4.5,-0.5)};
\draw  plot[smooth, tension=.7] coordinates {(2,4) (-0.5,1) (0.5,-2)};
\draw [dashed] (-0.5,5) -- (-2.5,-1) node[below]{Fenster};
\draw (-4.5,2.5) -- (2,0.5);
\draw  (-4.5,2.5) circle (0.2) node[above]{$\vec{w}_b$};
\draw  (2,0.5) circle (0.2) node[above]{$\vec{w}_s$};
\draw  (-1.5,0.5) circle (0.2) node[below]{$\vec{x}$};
\end{tikzpicture}
\end{center}
Falls $\vec{x}$ nicht in Fenster liegt, erfolgt \emph{keine} Adaption:
\begin{center}
\begin{tikzpicture}[scale=.5]
\draw  plot[smooth, tension=.7] coordinates {(-2.5,4.5) (-2.5,2) (-4.5,-0.5)};
\draw  plot[smooth, tension=.7] coordinates {(2,4) (-0.5,1) (0.5,-2)};
\draw [dashed] (-0.5,5) -- (-2.5,-1) node[below]{Fenster};
\draw (-4.5,2.5) -- (2,0.5);
\draw  (-4.5,2.5) circle (0.2) node[above]{$\vec{w}_b$};
\draw  (2,0.5) circle (0.2) node[above]{$\vec{w}_s$};
\draw  (-4,1) circle (0.2) node[left]{$\vec{x}$};
\end{tikzpicture}
\end{center}
$\vec{w}_b$ und $\vec{w}_s$ müssen verschiedene Klassen repräsentieren!
\end{itemize}
Fenster: Wird aus Relationen zwischen den Abständen $d_b=\Vert \vec{x}-\vec{w}_b\Vert$ und $d_s=\Vert \vec{x}-\vec{w}_s\Vert$ berechnet. Mit $0.2\leq v \leq 0.3$ muss $\min\left(\frac{d_b}{d_s}, \frac{d_s}{d_b}\right)>\frac{1-v}{1+v}$ sein, damit $\vec{x}$ im Fenster liegt. 
\subsubsection*{Lernalgorithmus}
\begin{enumerate}
\item Wähle ein $\vec{x}$.
\item Bestimme Best- und Second-Best-Matching-Neuronen.
\item Prüfe, ob $\vec{x}$ im Fenster liegt. Falls $\vec{x}$ im Fenster, dann:
\begin{anumerate}
\item $\operatorname{Klasse}(\vec{w}_b) = \operatorname{Klasse} (\vec{x}(t))$
$$\vec{w}_b(t+1) = \vec{w}_b(t)+\eta(t)\cdot \left[ \vec{x}(t) - \vec{w}_b(t)\right]$$
$$\vec{w}_s(t+1) = \vec{w}_s(t)-\eta(t)\cdot \left[ \vec{x}(t) - \vec{w}_s(t)\right]$$
\item $\operatorname{Klasse}(\vec{w}_s) = \operatorname{Klasse} (\vec{x}(t))$
$$\vec{w}_b(t+1) = \vec{w}_b(t)-\eta(t)\cdot \left[ \vec{x}(t) - \vec{w}_b(t)\right]$$
$$\vec{w}_s(t+1) = \vec{w}_s(t)+\eta(t)\cdot \left[ \vec{x}(t) - \vec{w}_s(t)\right]$$
\end{anumerate}
\end{enumerate}
LVQ 2.1 adaptiert ausschließlich die Referenzvektoren, die entlang der Klassengrenzen liegen.\\
$\To$ Die Referenzvektoren, die „innerhalb“ der Klassengebiete liegen, bleiben unverändert. $\To$ Gegebenenfalls schlechte Generalisierung innerhalb der Klassengebiete.

\subsection{LVQ 3}
Arbeitet analog zu LVG 2.1 und lässt zusätzlich den Fall zu, dass $\operatorname{Klasse}(\vec{w}_b) = \operatorname{Klasse}(\vec{w}_s) = \operatorname{Klasse} (\vec{x}(t))$ gilt. Für diesen Fall werden $\vec{w}_b$ und $\vec{w}_s$ adaptiert mit:
$$\vec{w}_{b/s}(t+1)=\vec{w}_{b/s}(t)+\eta(t) \cdot \epsilon \cdot \left[\vec{x}(t)-\vec{w}_{b/s}(t)\right]$$
Dabei ist $\epsilon \approx 0.5$.\\
Damit kann auch innerhalb der Klassen adaptiert werden. Somit ist die Repräsentation der Statistik der Klassengebiete möglich.

\subsection{Optimized LVQ (OLVQ)}
Ziel ist, dass jedes Datum der Trainingsdaten den gleichen Einfluss auf die Lage des Referenzvektors hat. Dazu führt man für jeden Referenzvektor eine eigene Schrittweite, die dies sicherstellt.

\subsection{Handhabung der LVQ-Verfahren}
Die Performance/Güte eines LVQ-Netzwerks hängt von folgenden Faktoren ab:
\begin{anumerate}
\item Zahl der Referenzvektoren pro Klasse 
\item Initialisierung der Referenzvektoren
\item Verwendeten LVQ-Verfahren
\item Zeitlichem Verlauf der Lernrate
\item Einem geeigneten Abbruchkriterium
\end{anumerate}

Als „best praxis“ Vorgehen zeit sich oft:
\begin{enumerate}
\item Mit OLVQ-Training beginnen.
\item Mit LVQ 1 bzw. LVQ 3 nachtrainieren.
\end{enumerate}

\chapter{Unüberwachte Lernverfahren}
\section{Abgrenzung überwachte Lernverfahren}
\begin{itemize}
\item MLP/RBF $\to$ supervised learning\\
$\vec{y}$ und $\vec{t}$ (Teacher) steuern das Neuronale Netzwerk.
\item LVQ $\to$ semi-supervised learning\\
Nutzen zwar Klassenzugehörigkeit der Beispieldaten aber geben keinen Teacher vor.
\item Unüberwachte Lernverfahren: unsupervised learning
\begin{itemize}
\item Ausschließlich auf Basis der Eingabedaten (absolut kein Teacher mehr).
\item \emph{Kompetitives Lernverfahren}: Neuronen konkurrieren um die Teilnahme am Adaptionsprozess.
\item Vektorquantisierung.
\end{itemize}
\end{itemize}

\section{Neural-Gas-Netzwerk (NG)}
Gas $\to$ Bezüge aus Thermodynamik und Molekülbewegungen.\\
Ziel: Anhand von $P$ Eingabedaten eine Menge von $M$ Referenzvektoren so verteilen, dass ein optimaler Vektorquantisierer entsteht (typischer Weise $P \gg M$).

\subsection{Optimale Vektorquantisierer}
Zum Bestimmen eines optimalen Vektorquantisierers wird ein Fehlermaß benötigt (Funktion). Der optimale Vektorquantisierer hat dann ein minimales Fehlermaß.\\
Quantisierungsfehler $E = \sum_{p=1}^P \| \vec{x}^p-\vec{w}_b(\vec{x}^p)\|$ (Abstand der Punkte vom Best-Matching-Neuron)\\
Man kann nachweisen, dass der NG-Algorithmus Optimalität sicherstellt.

\subsection{Algorithmus}
\lecdate{23.01.2018}
\begin{enumerate}
\item Wähle ein $\vec{x}$ aus der Menge der (gleichverteilten) Eingangsdaten.
\item Sortiere alle Referenzvektoren ihres euklidischen Abstandes zu $\vec{x}$.\\
$\To$ Ranking aller Referenzvektoren
\end{enumerate}
Ranking-Faktor: $k_i$, Lernreichweite: $h_\lambda(k_i)$.\\
Reichweite wird mit zunehmender Trainingsdauer kleiner. Das Best-Matching-Neuron hat die größte Adaption, weitere Neuronen eine geringere. Die Reichweite beschreibt dann, welche Neuronen beim Training noch (wie stark) einbezogen werden.
$$\Delta \vec{w}_i = \eta (t) \cdot h_\lambda(k_i) (\vec{x}-\vec{w}_i)$$\\
$$\lambda(t)=\lambda_{A}\cdot \left(\frac{\lambda_{E}}{\lambda_{A}}\right)^{\frac{t}{t_{max}}}$$
$$h_\lambda(k_i)=e^{\frac{-k_i}{\lambda(t)}}$$
$$\eta(t)=\eta_{A}\cdot \left(\frac{\eta_{E}}{\eta_{A}}\right)^{\frac{t}{t_{max}}}$$
Mit $\bullet_A$: Initialwert (Anfang), $\bullet_E$: Endwert, $t$: Trainingsepoche\\
Anmerkung: auf den Folien $\eta = \epsilon$.

\section{Kohonen-Netzwerke}
\begin{itemize}
\item Vektorquantisierer
\item benannt noch „Erfinder“
\item nutzt ncahbarschaft innerhalb der Netzwerk-Topologie
\item trainiert wird jeweils das Best-Matching-Neuron $\vec{r}'$ (Koordinate im Netzwerk) mit
$$\vec{r}'=\underset{\vec{r}}{\mathrm{argmin}}(\vec{x}-\vec{w}_{\vec{r}})$$
und dessen Nachbarn in der \emph{Topologie} $\To$ „Dimension“ des Netzwerks
\end{itemize}
Beispiel „Dimension“/Topologie: 
\begin{itemize}
\item 1D $\to$ Kette (von Neuronen)
\item 1,5D $\to$ Ring
\item 2D $\to$ Karte/Grid/Gitter
\end{itemize}
Definition Nachbarschaftsfunktionen:
$$h_{\vec{r}\vec{r}'}(t) = \exp (-\Vert\vec{r}'-\vec{r}\Vert)$$
Reichweite nimmt mit zunehmender Trainingsdauer ab.\\
$r(t)=r(t-1)\cdot \alpha$ mit $0<\alpha<1$, z.B. $\alpha=0,99$\\
$r(t=0)=10$ (als Richtwert)\\
$\eta(t)$: Analog zu Neural-Gas-Network oder einfacher $\eta(t)=\eta(0-1)\cdot \alpha$\\
$\eta(t=0)=0,8$ als Richtwert
$$\Delta \vec{w}_{\vec{r}} = \eta(t) \cdot h_{\vec{r}\vec{r}'}(t)\cdot (\vec{x}-\vec{w}_r)$$
Definition der Nachbarschaft (Dimension der Topologie) führt zur Eigenschaft der topologieerhaltenden Abbildung (bzw.: Nachbarschaften im Eingaberaum werden auf Nachbarschaften in der neuronalen Repräsentation abgebildet -- benachbarte Punkte bleiben benachbart).\\
Dies setzt voraus, dass die (Unter-)Mannigfaltigkeit (Dimensionen) der Eingabedaten mit der Dimension der Netzwerk-Topologie übereinstimmt. Sonst topologische Defekte zu erwarten (topologischer Defekt: Nachbarschaft in der neuronalen Repräsentation repräsentiert keine Nachbarschaft im Eingaberaum, bzw. umgekehrt). Beispiel Defekt: Kette (1D) in einer 2-D Ebene. Dort ist aber Beginn und Ende der Kette ggf. Nebeneinander, was nicht dem tatsächlichen Abstand entspricht.

Das Kohonen-Modell hat den stärksten biologischen Background.

\section{Growing-Neural-Gas}
$\To$ Netzwerk bestimmt selbst:
\begin{itemize}
\item Lage der Referenzvektoren
\item Anzahl der Referenzvektoren
\end{itemize}







\end{document}